% ============================================================
%  Hybrid PQ-OPAQUE — Scientific Paper
%  Template: Springer LNCS (llncs)
%  Compiler:  pdfLaTeX (Overleaf default)
% ============================================================
\documentclass{llncs}

% ── Language & Encoding (pdfLaTeX) ───────────────────────────
\usepackage[T2A]{fontenc}
\usepackage[utf8]{inputenc}
\usepackage[english,ukrainian]{babel}

% ── Mathematics ──────────────────────────────────────────────
\usepackage{amsmath}
\usepackage{amssymb}
\usepackage{mathtools}

% ── Tables ───────────────────────────────────────────────────
\usepackage{graphicx}
\usepackage{booktabs}
\usepackage{tabularx}
\usepackage{multirow}
\usepackage{array}
\usepackage{longtable}
\usepackage{makecell}
\usepackage{lscape}
\usepackage{tikz}
\usepackage{pgfplots}
\pgfplotsset{compat=1.18}
\usetikzlibrary{arrows.meta,positioning,shapes.geometric}

% ── Code Listings ────────────────────────────────────────────
\usepackage{listings}
\usepackage{xcolor}
\definecolor{eclblue}{RGB}{27,74,139}
\definecolor{eclgold}{RGB}{217,136,29}
\definecolor{eclgreen}{RGB}{49,125,92}
\definecolor{eclred}{RGB}{171,67,53}
\lstset{
  basicstyle=\ttfamily\footnotesize,
  breaklines=true,
  frame=single,
  framesep=3pt,
  backgroundcolor=\color{gray!8},
  numbers=none,
  keepspaces=true,
}

% ── Layout & Typography ──────────────────────────────────────
\usepackage{microtype}
\usepackage{caption}
\captionsetup{font=small, labelfont=bf}

% ── Hyperlinks ───────────────────────────────────────────────
\usepackage{hyperref}
\hypersetup{
  colorlinks=true,
  linkcolor=black,
  citecolor=black,
  urlcolor=blue,
}

% ── Custom column types ──────────────────────────────────────
\newcolumntype{L}[1]{>{\raggedright\arraybackslash}p{#1}}
\newcolumntype{C}[1]{>{\centering\arraybackslash}p{#1}}
\newcolumntype{R}[1]{>{\raggedleft\arraybackslash}p{#1}}

% ── Re-define theorem environments in Ukrainian ──────────────
% (llncs already provides theorem, definition, lemma)
\spnewtheorem*{sketchproof}{Доведення (ескіз)}{\itshape}{\rmfamily}

% ============================================================
\begin{document}
% ============================================================

\title{Гібридний постквантовий протокол автентифікації на основі
парольного ключового обміну OPAQUE з інтеграцією ML-KEM-768:
побудова, формальний аналіз та відтворювана реалізація}

\author{Олександр Мельниченко}

\institute{Ecliptix Security, Україна}

\maketitle

% ============================================================
\begin{abstract}
% ============================================================
Парольна автентифікація залишається одним із базових механізмів
підтвердження особи в сучасних інформаційних системах, однак
протоколи парольного автентифікованого ключового обміну (PAKE),
зокрема OPAQUE, спираються на класичні припущення складності для
Діффі--Геллмана та еліптичних кривих. Поява масштабованих квантових
обчислень робить ці припущення стратегічно вразливими, а сценарій
«збирай зараз~--- дешифруй пізніше» переносить проблему з
теоретичної площини у практичну.

У статті запропоновано Hybrid PQ-OPAQUE~--- реалізаційно орієнтоване
гібридне розширення OPAQUE, у якому ML-KEM-768 (FIPS~203) інтегровано
до 4DH-потоку автентифікованого ключового обміну. Сеансовий ключ
формується за допомогою HKDF-Extract із конкатенації чотирьох значень
Діффі--Геллмана на Ristretto255 і спільного секрету ML-KEM-768 із
транскриптно-залежною сіллю. На відміну від суто концептуальних описів
гібридних PAKE, у роботі зафіксовано повний формат подання
повідомлень, включно з 1-байтовим префіксом версії, контекстним
зв'язуванням \texttt{account\_id}, детермінованим виведенням
OPRF-ключа для конкретного облікового запису з глобального
\texttt{oprf\_seed} та окремим включенням постквантових елементів до
транскрипту і схеми виведення ключів.

Аргументація ґрунтується на трьох взаємодоповнювальних рівнях.
ProVerif~2.05 застосовано в двох окремих моделях для секретності та
автентифікації; Tamarin~Prover~1.10.0 використано для верифікованої
сурогатної 4DH-моделі, де архівований запуск підтверджує 8/8 лем за
22,83~с; Rust-реалізацію перевірено детерміністичними,
інтеграційними та властивісними тестами. Сукупно ці матеріали
формують сильну, хоча й не абсолютну, доказову базу: вони не
замінюють повного формального доведення коду реалізації, а твердження
про стійкість до офлайн-перебирання інтерпретуються лише для сценарію
компрометації бази даних без компрометації \texttt{oprf\_seed}.

Реалізацію виконано мовою \textbf{Rust~1.93.1} у вигляді
чотирикрейтного простору проєктів обсягом понад 5,3~тис. рядків
основного коду, майже 4,0~тис. рядків тестів та 1,46~тис. рядків
еталонних вимірювань. На платформі Apple~M1~Pro повний цикл
автентифікації становить 481,43~мс, повний раунд ML-KEM-768~---
122,01~мкс, тоді як домінантний внесок у латентність дає Argon2id
(510,86~мс в ізольованому вимірюванні). Постквантовий шар додає
159,3~мкс процесорного часу до повного ключового обміну без KSF та
збільшує повний обсяг рукостискання у форматі передавання до
2715~байт без зміни кількості раундів. Отримані результати
засвідчують, що перехід до постквантово-підсиленої парольної
автентифікації є практично досяжним уже нині, причому головний
компроміс пов'язаний насамперед із мережевим обсягом повідомлень, а не
з додатковими обчисленнями.
\end{abstract}

\keywords{OPAQUE, PAKE, ML-KEM-768, постквантова криптографія,
гібридний ключовий обмін, Ristretto255, HKDF, Tamarin, ProVerif, Rust}

% ============================================================
\section{Вступ}
% ============================================================

\subsection{Актуальність проблеми}

Парольна автентифікація залишається поширеним механізмом
підтвердження особи в сучасних розподілених інформаційних системах.
За різними оцінками, понад 80\% процедур автентифікації у веб-додатках
та мережевих сервісах базуються на паролях~\cite{bonneau2012,florencio2007}.
Протоколи парольного автентифікованого ключового обміну (PAKE)
забезпечують встановлення захищеного каналу між двома сторонами, що
поділяють спільний секрет низької ентропії~--- пароль~--- без розкриття
цього паролю противнику~\cite{bellare2000}. Серед відомих протоколів
PAKE протокол OPAQUE~\cite{jarecki2018}, запропонований Jarecki,
Krawczyk та Xu у 2018~році, належить до добре проаналізованих
представників класу асиметричних PAKE (у міжнародній літературі ---
augmented PAKE): сервер зберігає лише криптографічний запис, який не є
парольним еквівалентом при
компрометації лише бази даних~\cite{ietf-opaque}. Для практичних
реалізацій ця властивість потребує явного окреслення межі довіри:
компрометація глобального секрету, з якого деривуються OPRF-ключі,
відрізняється від компрометації пасивного сховища записів і не може
ототожнюватися з тим самим сценарієм безпеки.

\subsection{Квантова загроза}

У 1994~році Пітер Шор запропонував квантовий алгоритм~\cite{shor1994},
здатний розв'язувати задачу дискретного логарифмування за
поліноміальний час на квантовому комп'ютері. Цей алгоритм безпосередньо
загрожує всім криптосистемам на основі DLP та ECDLP:~RSA, ECDH, а
також будь-якому PAKE-протоколу, що будується на їх
основі~\cite{bernstein2017}. За оцінками, для зламу
Ristretto255/Curve25519 достатньо криптографічно значущого квантового
комп'ютера з приблизно 2\,330 логічними кубітами~\cite{roetteler2017}.
Водночас новіші оцінки для 256-бітної еліптичнокривої криптографії над
\texttt{secp256k1} вказують на ще жорсткіший сценарій: Google Quantum AI
у whitepaper 2026~року наводить дві реалізації алгоритму Шора для
ECDLP-256 з оцінками не більше ніж 1200~логічних кубітів і
90~млн Toffoli-операцій або не більше ніж 1450~логічних кубітів і
70~млн Toffoli-операцій~\cite{babbush2026}. Хоча ці числа не є прямою
оцінкою саме для Ristretto255, вони підсилюють загальний висновок:
криптосистеми на 256-бітних еліптичних кривих можуть виявитися
практично вразливими раніше, ніж це випливало з попереднього покоління
оцінок квантових ресурсів.

Особливу небезпеку становить стратегія «збирай зараз~--- дешифруй
пізніше»~\cite{mosca2018}: противник зберігає зашифрований трафік
сьогодні з метою ретроспективного криптоаналізу на квантовому
комп'ютері в майбутньому. Для систем, що опрацьовують довготривалі
конфіденційні дані, ця загроза є практично значущою вже зараз.

\subsection{Стандартизація постквантової криптографії}

У серпні 2024~року NIST опублікував перші стандарти постквантової
криптографії: FIPS~203 (ML-KEM)~\cite{fips203} та FIPS~204
(ML-DSA)~\cite{fips204}. ML-KEM базується на задачі навчання з
помилками на модульних решітках (Module-LWE), для якої не відомо
ефективних квантових алгоритмів. Гібридні підходи (класика~+ постквантова складова)
впроваджуються в TLS~1.3~\cite{stebila2024}, Signal
(PQXDH)~\cite{brendel2020} та WireGuard~\cite{hulsing2021}.

\subsection{Обмеження існуючих рішень}

Класичний OPAQUE у поточній специфікації IETF~\cite{ietf-opaque} не
містить механізмів захисту від квантових атак: ключовий матеріал
класичного 4DH може бути відновлений алгоритмом Шора. Інші
PAKE-протоколи (SRP, SPAKE2) є так само вразливими. Чисто постквантові
PAKE вже активно розвиваються: з'явилися компактні KEM-базовані
конструкції~\cite{arriaga2024}, UC/QROM-схеми для постквантового
середовища~\cite{lyu2024uc} та перші формалізовані підходи до
гібридного PAKE~\cite{hesse2024,lyu2024hybrid}. Водночас ці результати
ще не привели до усталеного стандарту асиметричного PAKE на кшталт OPAQUE.
Існуючі гібридні рішення (TLS~1.3, Signal, WireGuard) не адресують
парольну автентифікацію. Окремою проблемою є розрив між абстрактними
описами гібридних PAKE та відтворюваними матеріалами, придатними до
незалежної перевірки: у відкритій літературі бракує робіт, які одночасно
задають повний формат повідомлень протоколу, межі інтерпретації формальних
результатів, експериментальну базу та відкриту реалізацію.

\subsection{Мета та завдання дослідження}

\textbf{Мета:} розробка, формальний аналіз та реалізація гібридного
постквантового протоколу на основі OPAQUE, що забезпечує одночасну
стійкість до класичних та квантових атак без втрати фундаментальних
властивостей безпеки aPAKE.

\textbf{Завдання:}
\begin{enumerate}
  \item Розробити конструкцію гібридного комбінатора ключового матеріалу
        (4DH~+~ML-KEM-768).
  \item Специфікувати розширений транскрипт із постквантовими
        елементами, контекстним зв'язуванням \texttt{account\_id} та
        версіонуванням повідомлень протоколу.
  \item Побудувати багаторівневу доказову базу безпеки
        (редукційні ескізи, ProVerif, Tamarin, Rust-тести) з явним
        визначенням меж інтерпретації.
  \item Реалізувати протокол мовою Rust та підготувати відкритий набір
        дослідницьких матеріалів для відтворення результатів.
  \item Виміряти латентність, пропускну здатність та постквантові накладні
        витрати за допомогою статистичних еталонних вимірювань
        Criterion.
  \item Оцінити практичну придатність конструкції для інтерактивної
        парольної автентифікації на клієнтських і серверних вузлах.
\end{enumerate}

\subsection{Дослідницькі питання та критерії оцінювання}

Для систематизації роботи сформульовано чотири дослідницькі питання.
\begin{enumerate}
  \item[\textbf{RQ1}] Чи можна інтегрувати ML-KEM-768 до OPAQUE без
        збільшення кількості протокольних повідомлень і без втрати
        ключових властивостей aPAKE?
  \item[\textbf{RQ2}] Яким чином слід поєднувати класичний 4DH-матеріал і
        KEM-секрет так, щоб уникнути атаки зниження рівня безпеки та
        зберегти транскриптне зв'язування?
  \item[\textbf{RQ3}] Який обсяг формальної та експериментальної
        доказовості достатній, щоб робити стримані висновки про безпеку
        реалізації, не ототожнюючи символьну модель із кодом реалізації?
  \item[\textbf{RQ4}] Чи є постквантовий шар практично прийнятним для
        інтерактивної парольної автентифікації, якщо домінуючим
        компонентом латентності залишається Argon2id?
\end{enumerate}

Критеріями оцінювання є: (1) відсутність додаткових раундів обміну;
(2) наявність узгодженого опису мережевого формату повідомлень і криптографічного
транскрипту; (3) наявність формальних моделей і реалізаційних матеріалів, які
можна незалежно відтворити; (4) процесорні накладні витрати постквантового шару, що не
домінують над вартістю функції жорсткого перетворення пароля; (5)
чітко задокументована межа безпекових тверджень для сценаріїв
компрометації секретів сервера.

\subsection{Наукова новизна та внесок}

\begin{enumerate}
  \item \textbf{Гібридна 4DH~+~ML-KEM-768 конструкція.} Запропоновано
        реалізаційно орієнтований варіант Hybrid PQ-OPAQUE, у якому
        ML-KEM-768 є \emph{обов'язковою} складовою ключового обміну, а не
        факультативним розширенням.
  \item \textbf{Спеціалізований комбінатор ключового матеріалу.}
        $\mathrm{PRK} = \mathrm{HKDF\text{-}Extract}(\mathrm{salt}_\tau,\;
        dh_1 \| dh_2 \| dh_3 \| dh_4 \| ss_{\mathrm{kem}})$ із
        контекстно-залежною сіллю, що містить мічений геш розширеного
        транскрипту~$\tau$.
  \item \textbf{Специфікація рівня реалізації.} Окрім постквантових елементів,
        транскрипт включає окремий хеш контексту облікового запису,
        а мережевий формат визначає 1-байтовий префікс версії та точні
        розміри повідомлень.
  \item \textbf{Багаторівнева доказова база.} Для однієї й тієї самої
        конструкції узгоджено редукційні ескізи, дві ProVerif-моделі,
        Tamarin-модель із сурогатним поданням 4DH та Rust-тести.
  \item \textbf{Відкрита реалізація та пакет дослідницьких матеріалів.} Оновлений простір проєктів
        містить понад 5,3~тис. рядків основного Rust-коду, майже
        4,0~тис. рядків тестів і окремі модулі мікровимірювань,
        протокольних вимірювань та вимірювань пропускної здатності.
  \item \textbf{Репродуковані емпіричні результати.} Показано, що
        постквантовий шар збільшує процесорні витрати ключового обміну на
        159,3~мкс без KSF, тоді як домінуючим фактором латентності
        залишається Argon2id.
\end{enumerate}

\subsection{Термінологічні домовленості}

Щоб уникнути змішування англомовних і україномовних позначень, у
подальшому викладі використовується термінологія, зафіксована в
Табл.~\ref{tab:terminology}. Повторні вживання цих термінів слід
читати саме в наведених тут значеннях.

\begin{table}[htbp]
\caption{Термінологічні домовленості, прийняті в статті}
\label{tab:terminology}
\centering\small
\begin{tabularx}{\linewidth}{L{3.2cm}X}
\toprule
Термін & Зафіксоване значення \\
\midrule
Мережевий формат повідомлень & Повне байтове подання повідомлення на каналі, включаючи префікс версії та службові поля протоколу \\
Корисні дані & Вміст повідомлення без службового префікса версії та інших зовнішніх обгорток \\
Асиметричний PAKE (aPAKE) & Клас парольного автентифікованого ключового обміну, у якому сервер зберігає верифікатор, а не пароль у відкритому вигляді чи його прямий еквівалент \\
Схема виведення ключів & Послідовність процедур екстракції та розширення ключового матеріалу, що формує сеансові та допоміжні ключі \\
Пакет дослідницьких матеріалів & Сукупність вихідного коду, формальних моделей, тестів і еталонних вимірювань, достатніх для незалежного відтворення результатів \\
\bottomrule
\end{tabularx}
\end{table}

% ============================================================
\section{Аналіз літературних джерел та існуючих рішень}
% ============================================================

\subsection{Протоколи парольного автентифікованого ключового обміну}

Парольний автентифікований ключовий обмін (PAKE)~--- клас протоколів,
які дозволяють двом сторонам, що поділяють спільний секрет низької
ентропії (пароль), встановити автентифікований сесійний ключ, стійкий
до офлайн-атак за словником~\cite{boyko2000}. Протоколи PAKE
поділяються на дві категорії:

\begin{itemize}
  \item \emph{Збалансовані (balanced) PAKE}~--- обидві сторони зберігають
        ідентичне парольне представлення. Приклади: EKE~\cite{bellovin1992},
        SPAKE2~\cite{ladd2023}, CPace~\cite{haase2019}.
  \item \emph{Асиметричні PAKE (aPAKE; у міжнародній літературі --- augmented PAKE)}~--- сервер зберігає
        лише парольний верифікатор. Приклади: SRP~\cite{wu1998},
        AuCPace~\cite{haase2019}, OPAQUE~\cite{jarecki2018}.
\end{itemize}

\begin{table}[htbp]
\caption{Порівняльна характеристика протоколів PAKE}
\label{tab:pake-comparison}
\centering\small
\begin{tabularx}{\linewidth}{L{3.5cm}CCCCC}
\toprule
Властивість & EKE & SRP & SPAKE2 & OPAQUE \\
\midrule
Тип & Збалансований & Асиметричний & Збалансований & Асиметричний \\
Нееквівалентний серверний запис & Ні & Частково & Ні & Так \\
Стійкість до попередніх обчислень & Ні & Ні & Ні & Так \\
Пряма секретність & Так & Так & Так & Так \\
Стійкість до компрометації серверного сховища & Ні & Часткова & Ні & У моделі пасивної БД \\
Формальне доведення безпеки & IC & Відсутнє & ROM & UC \\
Квантова стійкість & Ні & Ні & Ні & Ні \\
\bottomrule
\end{tabularx}
\end{table}

У Табл.~\ref{tab:pake-comparison} рядок для OPAQUE щодо компрометації
серверного сховища слід читати у вузькому, але стандартному сенсі
компрометації лише бази записів автентифікації. Він не охоплює
компрометацію глобального серверного секрету, з якого можуть
деривуватися OPRF-ключі.

\subsection{Протокол OPAQUE: формальний опис}

Протокол OPAQUE~\cite{jarecki2018} складається з фази реєстрації та
фази автентифікації, обидві спираються на механізм забутливої
псевдовипадкової функції (OPRF). Стандартний IETF OPAQUE використовує потрійний обмін Діффі--Геллмана (3DH):
$dh_1 = sk_S \cdot PK_C$ (взаємна автентифікація), $dh_2 = sk_S \cdot E_C$,
$dh_3 = e_S \cdot PK_C$ (пряма секретність). У даній роботі додано
\textbf{четвертий добуток} $dh_4 = e_S \cdot E_C$, який додатково
зв'язує обидві ефемерні компоненти в одному транскрипті і на рівні
4DH-інтуїції покликаний ускладнити UKS-подібні підстановки; у цій
роботі це тлумачиться як дизайн-мотивація, а не як окремо доведена
властивість. Це розширення є внеском даного дослідження.
Безпека базового OPAQUE доведена в моделі Universal Composability
(UC)~\cite{canetti2001}.

\textbf{Забутлива псевдовипадкова функція (OPRF).} Реалізація
використовує гешування в групу Ristretto255~\cite{devalence2024}:
\begin{equation}
  \alpha = r \cdot H'(\mathrm{ctx} \| x),\quad
  \beta = k \cdot \alpha,\quad
  F_k(x) = r^{-1} \cdot \beta,
\end{equation}
де $r \xleftarrow{\$} \mathbb{Z}_q^*$~--- засліплюючий скаляр
клієнта, $k$~--- секретний ключ сервера.

\subsection{Постквантова криптографія та стандарти NIST 2024}

ML-KEM-768 обрано на рівні безпеки NIST рівня~3 (еквівалент AES-192).
Параметри: $|pk| = 1184$~байт, $|ct| = 1088$~байт, $|ss| = 32$~байти,
безпека IND-CCA2. Конструкція:
\begin{equation}
  \mathrm{KeyGen}() \to (pk, sk);\quad
  \mathrm{Encaps}(pk) \to (ct, ss);\quad
  \mathrm{Decaps}(sk, ct) \to ss.
\end{equation}

\subsection{Обґрунтування вибору ML-KEM-768}

Вибір ML-KEM-768 у цій роботі не є випадковим. З одного боку,
ML-KEM-512 пропонує менші повідомлення і нижчі накладні витрати, але
орієнтується на NIST рівень~1. З іншого боку, ML-KEM-1024 надає вищий
рівень запасу безпеки ціною відчутно більшого розміру повідомлень у мережевому форматі
і більшого навантаження на транспорт. Для інтерактивної парольної
автентифікації, де протокол уже містить відносно дорогий KSF,
центральним стає не максимальний можливий запас безпеки сам по собі, а
збалансованість між криптографічною стійкістю, розміром повідомлень та
операційною придатністю.

За сукупністю цих міркувань ML-KEM-768 постає як виважений
вибір для Hybrid PQ-OPAQUE. Цей параметричний варіант посідає
проміжне місце в сімействі ML-KEM і надає дві важливі
переваги. По-перше, рівень NIST~3 добре узгоджується з практикою
довготривалого захисту облікових записів і секретів сесії. По-друге,
розміри $pk_{\mathrm{kem}}$ та $ct_{\mathrm{kem}}$ хоча й істотно
збільшують обсяг повідомлень порівняно з класичним OPAQUE, проте не
роблять протокол непридатним для інтерактивного застосування.
Таким чином, ML-KEM-768 у цьому дизайні є не ``максимально
потужним'' і не ``мінімально легким'' варіантом, а параметром із
збалансованим співвідношенням запасу безпеки та системних накладних витрат.

\subsection{Нові підходи до постквантових PAKE (2023--2025)}

У 2023--2025~роках у цій галузі виокремилося кілька головних
дослідницьких ліній.

\textbf{KEM-базовані постквантові PAKE.} Одна з найактивніших ліній
робіт спирається на абстракцію KEM і прагне побудувати PAKE з
мінімальними накладними витратами щодо базового постквантового
примітива. Репрезентативним прикладом є CHIC~\cite{arriaga2024}, де автори
одержують компактний решітковий PAKE з доведенням UC-безпеки та
реалізаційною оцінкою. Такі результати демонструють, що чисто
постквантовий PAKE уже вийшов за межі суто теоретичних побудов, однак
типово не орієнтується на асиметричну клієнт-серверну модель OPAQUE.

\textbf{Постквантові UC/QROM-конструкції.} Окремий напрямок зосереджено
на досягненні сильних композиційних гарантій безпеки. Робота
\cite{lyu2024uc} пропонує UC-безпечні PAKE для постквантового світу,
включаючи QROM-варіанти на основі LWE. Ці результати важливі як
теоретична опора, але вони не розв'язують безпосередньо задачу
міграції вже існуючих OPAQUE-подібних систем, де ключовими є сумісність
з реальною серверною моделлю експлуатації, мережевий формат повідомлень
і відтворювані дослідницькі матеріали.

\textbf{Решітково-орієнтовані OPRF-замінники та KEM-to-PAKE компілятори.}
Конструкції на основі задачі Learning With Errors (LWE) пропонуються як
постквантові альтернативи Ristretto-OPRF, проте мають суттєво більший розмір
сліпих елементів та знижену ефективність~\cite{beguinet2023}. Водночас
останні роботи вказують, що узагальнені KEM-to-PAKE підходи потребують дуже
обережного поводження з ідеальним шифром, хешуванням у простір
публічних ключів та інстанціюванням низькоентропійних секретів~\cite{arriaga2024}.

\textbf{Ізогенні PAKE.} Атаки Castryck--Decru (2022)~\cite{castryck2023}
підірвали довіру до SIDH-базованих конструкцій. SIDH-базовані PAKE
вважаються небезпечними.

\textbf{Гібридний підхід як окремий напрямок досліджень.} Новітні праці
вже формалізують саму постановку гібридного PAKE. Зокрема, робота
\cite{hesse2024} пропонує PAKE-комбінатори з ефективними
постквантовими інстанціюваннями, а~\cite{lyu2024hybrid} аналізує
гібридний PAKE в UC-рамці через паралельну та послідовну композицію.
Показово, що у 2025~році з'явилася й індивідуальна чернетка IETF
із гібридного постквантового асиметричного PAKE для CFRG~\cite{vos2025},
що свідчить про перехід теми від суто академічних конструкцій до
передстандартизаційної фази. На цьому тлі Hybrid PQ-OPAQUE
позиціонується як реалізаційно-орієнтований крок у напрямі гібридного
розширення OPAQUE з явною специфікацією мережевого формату,
формальними моделями та відкритою реалізацією.

% ============================================================
\section{Постановка задачі}
% ============================================================

\subsection{Формулювання проблеми та модель загроз}

Нехай $\mathcal{C}$ (клієнт/ініціатор) та $\mathcal{S}$
(сервер/респондер)~--- дві сторони протоколу, де $\mathcal{C}$ володіє
паролем $\mathrm{pwd}$, а $\mathcal{S}$ зберігає реєстраційний запис
$\mathrm{rec}$. Позначимо через $\mathbb{G}$~--- групу Ristretto255
простого порядку~$q$ з генератором~$g$.

\textbf{Активи, що підлягають захисту:}
\begin{itemize}
  \item Пароль $\mathrm{pwd}$; статичні ключі $sk_C, sk_S$; ефемерні
        ключі $e_C, e_S, sk_{\mathrm{kem}}$;
  \item Реєстраційний запис $\mathrm{rec} = (\mathrm{envelope}, pk_C)$;
  \item Сесійний ключ $K_s$, майстер-ключ $K_m$, спільний секрет
        $ss_{\mathrm{kem}}$.
\end{itemize}

\textbf{Модель противника.} Розглядаємо три класи:
\begin{itemize}
  \item \emph{Класичний $\mathcal{A}_C$}: поліноміально обмежений,
        активний «людина посередині», може компрометувати серверне
        сховище;
  \item \emph{Квантовий $\mathcal{A}_Q$}: розв'язує DLP/ECDLP (алгоритм
        Шора), застосовує стратегію «збирай зараз~--- дешифруй пізніше»;
  \item \emph{Композитний $\mathcal{A}_{CQ}$}: активний «людина
        посередині» із квантовим обчислювальним оракулом.
\end{itemize}

\textbf{Вимоги до протоколу:}
\begin{description}
  \item[R1] (OPAQUE-властивості): таємність пароля, пряма секретність,
        взаємна автентифікація, стійкість до офлайн-перебирання;
  \item[R2] (Постквантова стійкість): захист від стратегії «збирай
        зараз~--- дешифруй пізніше»;
  \item[R3] (AND-інтерпретація): комбінатор ключового матеріалу має
        бути сумісним з AND-інтерпретацією в прийнятій моделі, тобто
        явне відновлення фінального ключа не повинно зводитися до
        ослаблення лише одного шару;
  \item[R4] (Мінімальні накладні витрати): збереження 3 повідомлень та
        1,5 раундів обміну;
  \item[R5] (Реалізованість): кросплатформна реалізація, придатна до
        відтворення та подальшого інженерного розвитку.
\end{description}

\subsection{Формальні визначення}

\begin{definition}[KEM]
$(\mathrm{KeyGen}, \mathrm{Encaps}, \mathrm{Decaps})$:
$\mathrm{KeyGen}() \to (pk, sk)$;
$\mathrm{Encaps}(pk) \to (ct, ss)$;
$\mathrm{Decaps}(sk, ct) \to ss$.
IND-CCA2 безпека: жоден поліноміально обмежений противник з оракулом
декапсуляції не відрізняє реальний $ss$ від випадкового.
\end{definition}

\begin{definition}[Гібридний aPAKE у межах цієї роботи]
Протокол $\Pi$ між $\mathcal{C}(\mathrm{pwd})$ та
$\mathcal{S}(\mathrm{rec})$, у якому
\[
  \mathrm{IKM}_{\mathrm{cl}} = dh_1 \| dh_2 \| dh_3 \| dh_4,\qquad
  \mathrm{IKM}_{\mathrm{pq}} = ss_{\mathrm{kem}},
\]
а фінальний пре-ключ формується як
\[
  \mathrm{PRK} = \mathrm{HKDF\text{-}Extract}\!\bigl(
    \mathrm{salt}_{\tau},\;
    \mathrm{IKM}_{\mathrm{cl}} \| \mathrm{IKM}_{\mathrm{pq}}\bigr).
\]
Подальший розділ аналізує, якою мірою така конструкція сумісна з
AND-інтерпретацією безпеки в прийнятій моделі.
\end{definition}

\begin{definition}[Псевдовипадковість HKDF-Extract]
$\mathrm{HKDF\text{-}Extract}(\mathrm{salt}, \mathrm{IKM})$ на основі
HMAC-SHA-512 моделюється як екстрактор: якщо $\mathrm{IKM}$ містить
приховану обчислювальну ентропію відносно противника, тоді вихід PRK є
обчислювально нерозрізнюваним від рівномірного за перевагою не
більше~$\varepsilon$.
\end{definition}

\textbf{Нотація:} $\mathbb{G}$~--- Ristretto255, $|\mathbb{G}| = q$;
малі літери~--- скаляри; великі~--- точки кривої; $\|$~---
конкатенація; $H(\cdot)$~--- SHA-512;
$\mathcal{K} = (K_s, K_m)$~--- сесійний та майстер-ключ.

% ============================================================
\section{Опис запропонованого протоколу Hybrid PQ-OPAQUE}
% ============================================================

\subsection{Архітектура та огляд}

Гібридний PQ-OPAQUE є гібридним постквантовим розширенням
OPAQUE~\cite{ietf-opaque}. Протокол зберігає властивість aPAKE (сервер
не отримує доступу до пароля) та додає стійкість до стратегії «збирай
зараз~--- дешифруй пізніше». Протокол має дві фази: \emph{реєстрація}
(одноразова фаза) та \emph{автентифікація} (обмін із трьох
повідомлень KE1--KE2--KE3). Надалі терміни «мережевий формат
повідомлень» і «корисні дані» вживаються у значеннях, зафіксованих у
Табл.~\ref{tab:terminology}.

\begin{table}[htbp]
\caption{Криптографічні примітиви протоколу Hybrid PQ-OPAQUE}
\label{tab:primitives}
\centering\small
\begin{tabularx}{\linewidth}{L{3.2cm}L{3.2cm}L{2.8cm}L{2cm}}
\toprule
Компонент & Примітив & Реалізація & Рівень безпеки \\
\midrule
Група OPRF & Ristretto255 & \texttt{curve25519-dalek} & 128 біт \\
Геш-функція & SHA-512 & \texttt{sha2} & 256 біт \\
KDF (витяг) & HKDF-Extract (HMAC-SHA-512) & \texttt{hmac + sha2} & 256 біт \\
KDF (розш.) & HKDF-Expand (HMAC-SHA-512) & \texttt{hmac + sha2} & 256 біт \\
MAC & HMAC-SHA-512 & \texttt{hmac + sha2} & 256 біт \\
KSF & Argon2id (MODERATE) & \texttt{argon2} & Адаптивний \\
AEAD & XSalsa20-Poly1305 & \texttt{crypto\_secretbox} & 256 біт \\
KEM & ML-KEM-768 (FIPS~203) & \texttt{ml-kem} & NIST Рівень~3 \\
Класичний обмін & 4DH (Ristretto255) & \texttt{curve25519-dalek} & 128 біт \\
Рівночасове порівняння & Побайтове без раннього виходу & \texttt{subtle} & --- \\
\bottomrule
\end{tabularx}
\end{table}

\begin{table}[htbp]
\caption{Розміри повідомлень автентифікації в корисних даних і мережевому форматі}
\label{tab:message-sizes}
\centering\small
\begin{tabular}{lrrrr}
\toprule
Повідомлення & Класичні корисні дані & Гібридні корисні дані & Гібридний мережевий формат & Постквантові накладні витрати \\
\midrule
KE1 & 88 байт & 1272 байти & 1273 байти & $+1184$ ($pk_{\mathrm{kem}}$) \\
KE2 & 288 байт & 1376 байтів & 1377 байтів & $+1088$ ($ct_{\mathrm{kem}}$) \\
KE3 & 64 байти & 64 байти & 65 байтів & $0$ \\
\midrule
\textbf{Разом} & \textbf{440 байтів} & \textbf{2712 байтів} & \textbf{2715 байтів} & \textbf{$+2272$ ($+516\%$)} \\
\bottomrule
\end{tabular}
\end{table}

\begin{figure}[htbp]
\centering
\begin{tikzpicture}
\begin{axis}[
  width=0.92\linewidth,
  height=6.0cm,
  ybar,
  bar width=10pt,
  ymin=0,
  ymax=1500,
  ylabel={Байти},
  symbolic x coords={KE1,KE2,KE3},
  xtick=data,
  enlarge x limits=0.22,
  ymajorgrids=true,
  grid style={gray!20},
  legend style={at={(0.5,1.02)}, anchor=south, legend columns=2, draw=none, font=\small},
  tick label style={font=\small},
  label style={font=\small}
]
\addplot[fill=eclblue!75, draw=eclblue] coordinates {
  (KE1,88) (KE2,288) (KE3,64)
};
\addplot[fill=eclgold!85, draw=eclgold!70!black] coordinates {
  (KE1,1273) (KE2,1377) (KE3,65)
};
\legend{Класичні корисні дані,Гібридний мережевий формат}
\end{axis}
\end{tikzpicture}
\caption{Порівняння обсягу повідомлень автентифікації: додатковий мережевий
обсяг зосереджено у повідомленнях KE1 та KE2, тоді як KE3 лишається
майже незмінним.}
\label{fig:wire-size-growth}
\end{figure}

Рисунок~\ref{fig:wire-size-growth} робить наочним важливий практичний
висновок: постквантовий шар не змінює суттєво логіку завершення
автентифікації, але істотно збільшує транспортний бюджет перших двох
повідомлень, тобто саме тих фаз, де відбувається передавання
$pk_{\mathrm{kem}}$ та $ct_{\mathrm{kem}}$.

\subsection{Фаза реєстрації}

\textbf{Крок~1 ($\mathcal{C} \to \mathcal{S}$).} Генерація статичної
пари $(sk_C, PK_C)$; OPRF-засліплення:
\begin{equation}
  r \xleftarrow{\$} \mathbb{Z}_q^*, \qquad
  B = r \cdot H'\!\bigl(\text{``ECLIPTIX-OPAQUE-v1/OPRF''} \| \mathtt{0x00} \| \mathrm{pwd}\bigr).
\end{equation}
$\mathrm{RegistrationRequest} = B$ (32~байти).

\textbf{Крок~2 ($\mathcal{S} \to \mathcal{C}$).} Виведення OPRF-ключа:
\begin{equation}
  s_{\mathrm{oprf}} =
  \mathrm{HMAC}\!\bigl(\mathrm{oprf\_seed},
    \text{``ECLIPTIX-OPAQUE-v1/OPRF-Seed''}\bigr)[0..31],
\end{equation}
\begin{equation}
  k_{\mathrm{oprf}} = \mathrm{ScalarReduce}\!\bigl(
    \mathrm{HMAC}\!\bigl(s_{\mathrm{oprf}},
    \text{``ECLIPTIX-OPAQUE-v1/OPRF-Key''} \|
    \mathrm{account\_id} \| \mathrm{ctr}\bigr)\bigr),
\end{equation}
де $\mathrm{ctr}$~--- найменший лічильник, для якого отриманий скаляр
ненульовий. Обчислення $Z = k_{\mathrm{oprf}} \cdot B$; ключ
$k_{\mathrm{oprf}}$ негайно знищується. Відповідь $= Z \| PK_S$
(64~байти).

\textbf{Крок~3 ($\mathcal{C}$).} OPRF-розсліплення
$U = r^{-1} \cdot Z$; OPRF-вихід:
\begin{equation}
  y_{\mathrm{oprf}} = H\!\bigl(\text{``ECLIPTIX-OPAQUE-v1/OPRF''} \|
  \mathtt{0x01} \| \mathrm{pwd} \| U\bigr).
\end{equation}
Вхід та сіль для KSF:
\begin{equation}
  x_{\mathrm{ksf}} = H\!\bigl(
    \text{``ECLIPTIX-OPAQUE-v1/KSF''} \|
    y_{\mathrm{oprf}} \| \mathrm{pwd}\bigr),
\end{equation}
\begin{equation}
  \mathrm{salt}_{\mathrm{ksf}} =
  H\!\bigl(\text{``ECLIPTIX-OPAQUE-v1/KSF-Salt''} \|
  y_{\mathrm{oprf}}\bigr)[0..15].
\end{equation}
Тоді рандомізований пароль обчислюється як
\begin{equation}
  \mathrm{rwd} = \mathrm{Argon2id}\!\bigl(
    x_{\mathrm{ksf}},\; \mathrm{salt}_{\mathrm{ksf}};\;
    m = 256~\mathrm{MiB},\; t = 3,\; p = 1\bigr).
\end{equation}
Конверт запечатується за допомогою XSalsa20-Poly1305 і ключа,
виведеного через HKDF;
$\mathrm{RegistrationRecord} = \mathrm{Envelope} \| PK_C$
(168~байт корисних даних, 169~байт у мережевому форматі з урахуванням префікса версії).

\subsection{Прив'язка \texttt{account\_id}, версії та контексту}

На рівні реалізації протокол фіксує додатковий контекст, який у
багатьох абстрактних описах OPAQUE лишається неявним. По-перше, всі повідомлення у мережевому форматі
мають 1-байтовий префікс версії $v = \mathtt{0x01}$, який відокремлює
корисні дані від повного формату передавання і забезпечує безпечне відхилення
невідомих версій. По-друге, протокол оперує явним ідентифікатором
облікового запису \texttt{account\_id}, який використовується не лише як
ключ доступу до запису користувача, а й як криптографічно зв'язаний
контекст.

Нехай
\begin{equation}
  \eta = H\!\bigl(\text{``ECLIPTIX-OPAQUE-v2/AccountContext''} \| \mathrm{account\_id}\bigr).
\end{equation}
Різні доменно-розділові мітки використано тут навмисно. Префікси
\texttt{ECLIPTIX-OPAQUE-v1/\dots} зарезервовано для базових контекстів
OPRF, KSF і транскрипту; \texttt{ECLIPTIX-OPAQUE-v2/AccountContext}
маркує окремо введену прив'язку \texttt{account\_id}, щоб не змішувати
її з ранніми варіантами серіалізації цього контексту; префікси
\texttt{ECLIPTIX-OPAQUE-PQ-v1/\dots} відокремлюють уже власне
гібридний комбінатор і похідні ключі. Тому номер версії в такій мітці
слід читати як локальний ідентифікатор домену, а не як глобальний номер
всієї схеми.
Тоді \texttt{account\_id} впливає на протокол у двох точках:
\begin{enumerate}
  \item через виведення OPRF-ключа, пов'язаного з конкретним обліковим записом,
        $k_{\mathrm{oprf}} = \mathrm{KDF}(\mathrm{oprf\_seed}, \mathrm{account\_id})$;
  \item через включення $\eta$ до транскрипту автентифікації, а отже і до
        похідних ключів $K_s$, $K_m$, $K_{\mathrm{mac}}^S$,
        $K_{\mathrm{mac}}^C$.
\end{enumerate}
Таке зв'язування унеможливлює перенесення елементів мережевого формату між різними
обліковими записами навіть за однакового пароля, а також робить опис
реалізації точнішим порівняно з абстрактними специфікаціями, які залишають
ідентифікатор користувача поза криптографічною частиною рукостискання.

\subsection{Фаза автентифікації та гібридний ключовий обмін}

\subsubsection{Генерація KE1 ($\mathcal{C} \to \mathcal{S}$)}

\begin{equation}
  e_C \xleftarrow{\$} \mathbb{Z}_q^*,\quad
  E_C = e_C \cdot g;\qquad
  (pk_{\mathrm{kem}}, sk_{\mathrm{kem}}) \leftarrow
    \mathrm{ML\text{-}KEM\text{-}768.KeyGen}();\qquad
  n_C \xleftarrow{\$} \{0,1\}^{192}.
\end{equation}
\begin{equation}
  \mathrm{KE1} = B \| E_C \| n_C \| pk_{\mathrm{kem}}
  \quad(32 + 32 + 24 + 1184 = 1272~\text{байти корисних даних}).
\end{equation}
Повідомлення у мережевому форматі має вигляд $v \| \mathrm{KE1}$ і займає 1273~байти.

\subsubsection{Генерація KE2 ($\mathcal{S} \to \mathcal{C}$)}

\textbf{Чотиристоронній Діффі--Геллман (4DH):}
\begin{equation}
  dh_1 = sk_S \cdot PK_C,\quad
  dh_2 = sk_S \cdot E_C,\quad
  dh_3 = e_S \cdot PK_C,\quad
  dh_4 = e_S \cdot E_C.
\end{equation}

\textbf{Інкапсуляція ML-KEM-768:}
\begin{equation}
  (ct_{\mathrm{kem}}, ss_{\mathrm{kem}})
    \leftarrow \mathrm{ML\text{-}KEM\text{-}768.Encaps}(pk_{\mathrm{kem}}).
\end{equation}

\textbf{Розширений транскрипт} (включає постквантові елементи):
\begin{align}
  \tau &= E_C \| E_S \| n_C \| n_S \| PK_C \| PK_S \|
         \mathrm{CredResp} \| pk_{\mathrm{kem}} \| ct_{\mathrm{kem}} \| \eta,\\
  \tau_H &= H\!\bigl(\text{``ECLIPTIX-OPAQUE-v1/Transcript''} \| \tau\bigr).
\end{align}

\textbf{Гібридний комбінатор:}
\begin{align}
  \mathrm{combined\_ikm} &= dh_1 \| dh_2 \| dh_3 \| dh_4 \| ss_{\mathrm{kem}}
  \quad(128 + 32 = 160~\text{байт}),\\
  \mathrm{salt} &= \text{``ECLIPTIX-OPAQUE-PQ-v1/Combiner''} \| \tau_H,\\
  \mathrm{PRK} &= \mathrm{HKDF\text{-}Extract}(\mathrm{salt},
    \mathrm{combined\_ikm}).
\end{align}

\textbf{Деривація ключів:}
\begin{align}
  K_s &= \mathrm{HKDF\text{-}Expand}(\mathrm{PRK}, \text{``SessionKey''}, 64),\\
  K_m &= \mathrm{HKDF\text{-}Expand}(\mathrm{PRK}, \text{``MasterKey''}, 32),\\
  K_{\mathrm{mac}}^S &= \mathrm{HKDF\text{-}Expand}(\mathrm{PRK},
    \text{``ResponderMAC''}, 64),\\
  K_{\mathrm{mac}}^C &= \mathrm{HKDF\text{-}Expand}(\mathrm{PRK},
    \text{``InitiatorMAC''}, 64).
\end{align}
\emph{Примітка.} Для стислості мітки подано в скороченому вигляді.
Повні доменно-розділові мітки: \texttt{ECLIPTIX-OPAQUE-PQ-v1/SessionKey},
\texttt{ECLIPTIX-OPAQUE-PQ-v1/MasterKey} тощо.
\begin{equation}
  \mathrm{KE2} = n_S \| E_S \| \mathrm{CredResp} \|
    \mathrm{HMAC}(K_{\mathrm{mac}}^S, \tau) \| ct_{\mathrm{kem}}
  \quad(1376~\text{байтів корисних даних}).
\end{equation}
Повідомлення у мережевому форматі має вигляд $v \| \mathrm{KE2}$ і займає 1377~байтів.

\subsubsection{Генерація KE3 ($\mathcal{C} \to \mathcal{S}$)}

\begin{enumerate}
  \item OPRF-фіналізація, відкриття конверта, відновлення
        $(PK_S, sk_C, PK_C)$;
  \item Дзеркальний 4DH: $dh_1 = sk_C \cdot PK_S$,
        $dh_2 = e_C \cdot PK_S$, $dh_3 = sk_C \cdot E_S$,
        $dh_4 = e_C \cdot E_S$ (додаткове зв'язування ефемерних
        компонент на рівні 4DH-інтуїції);
  \item Декапсуляція: $ss_{\mathrm{kem}} =
        \mathrm{ML\text{-}KEM\text{-}768.Decaps}(sk_{\mathrm{kem}},
        ct_{\mathrm{kem}})$; $sk_{\mathrm{kem}}$ негайно знищується;
  \item Ідентична деривація PRK; верифікація
        $\mathrm{HMAC}(K_{\mathrm{mac}}^S, \tau)$ рівночасовою
        функцією побайтового порівняння без раннього виходу;
  \item $\mathrm{KE3} = \mathrm{HMAC}(K_{\mathrm{mac}}^C, \tau)$
        (64~байти корисних даних, 65~байтів із префіксом версії).
\end{enumerate}

\subsubsection{Завершення на стороні сервера}

Верифікація KE3 рівночасовою функцією порівняння. При успіху~---
пара $(K_s, K_m)$ стає доступною обом сторонам, підтверджуючи
взаємну автентифікацію. При відмові~--- всі ключові матеріали
негайно знищуються через безпечне затирання.

\subsection{Включення постквантових елементів у транскрипт}

Включення $pk_{\mathrm{kem}}$ та $ct_{\mathrm{kem}}$ у транскрипт є
\textbf{криптографічно необхідним}: воно \emph{криптографічно зв'язує}
постквантовий обмін з класичним. Без цього зв'язування противник типу
«людина посередині» може замінити $pk_{\mathrm{kem}}$ власним ключем на
рівні ML-KEM. Окреме включення хешу контексту облікового запису $\eta$
додатково унеможливлює перенесення валідного транскрипту між різними
користувачами або повторне використання повідомлень автентифікації для
іншого \texttt{account\_id}. Отже, транскрипт у цій роботі виконує
дві взаємопов'язані функції: він зв'язує класичний і постквантовий шари між собою та одночасно
прив'язує їх до конкретного облікового запису й версії мережевого
формату.

% ============================================================
\section{Аналіз безпеки}
% ============================================================

\subsection{Формальний аналіз властивостей безпеки}

\paragraph{Сценарій офлайн-перебирання за фіксованої межі довіри.}
У прийнятій моделі загроз нехай противник отримав пасивне серверне
сховище $\mathrm{rec}$ і спостережувані транскрипти, але не
\texttt{oprf\_seed}. Тоді він не може перевіряти словникові кандидати
офлайн, окрім як із нехтовно малою перевагою, зумовленою порушенням
забутливості OPRF або стійкості конверта. Якщо ж \texttt{oprf\_seed}
скомпрометовано, задача зводиться до словникового пошуку з питомою
вартістю одного припущення, що домінується Argon2id. Це формулювання
має статус обмеженого безпекового твердження для чітко визначеної межі
довіри, а не універсального твердження для всіх реалізацій OPAQUE.

\begin{sketchproof}
\textbf{Випадок 1: \texttt{oprf\_seed} не скомпрометовано.} Без
OPRF-ключа, пов'язаного з конкретним обліковим записом, противник не може обчислити
$y_{\mathrm{oprf}}$, а отже і відтворити
$x_{\mathrm{ksf}}, \mathrm{salt}_{\mathrm{ksf}}$ та
$\mathrm{rwd}$. Тому запис конверта не дає ефективного офлайн-оракула
перевірки пароля. Заміна OPRF-виходу випадковим значенням дає
\[
  |\Pr[S_1] - \Pr[S_0]| \leq
  \mathrm{Adv}_{\mathrm{OPRF}}^{\mathrm{obliv}}
  \leq \mathrm{Adv}_{\mathrm{Ristretto255}}^{\mathrm{CDH}}.
\]
\textbf{Гра~1 $\to$ Гра~2:} замінюємо конверт шифруванням
випадкового рядка (IND-CPA XSalsa20-Poly1305):
\[
  |\Pr[S_2] - \Pr[S_1]| \leq
  \mathrm{Adv}_{\mathrm{XSalsa20}}^{\mathrm{IND\text{-}CPA}}.
\]
\textbf{Наслідок.} За відсутності компрометації
\texttt{oprf\_seed} найкраща стратегія зводиться до онлайн-вгадувань,
тому
\[
  \mathrm{Adv}^{\mathrm{dict}} \leq
  \mathrm{Adv}_{\mathrm{Ristretto255}}^{\mathrm{CDH}} +
  \mathrm{Adv}_{\mathrm{XSalsa20}}^{\mathrm{IND\text{-}CPA}} +
  \frac{q_{\mathrm{on}}}{|\mathcal{D}|}.
\]
\textbf{Випадок 2: \texttt{oprf\_seed} скомпрометовано.} Тоді кожне
словникове припущення $\mathrm{pwd}'$ може бути перевірене офлайн через
обчислення $y_{\mathrm{oprf}}'$, далі $x_{\mathrm{ksf}}'$,
$\mathrm{salt}_{\mathrm{ksf}}'$ і одну спробу відкриття конверта. Отже
вартість одного припущення асимптотично та практично домінується
Argon2id:
\[
  T_{\mathrm{guess}} \approx T_{\mathrm{Argon2id}},\qquad
  C_{\mathrm{offline}} \approx |\mathcal{D}| \cdot T_{\mathrm{guess}}.
  \qquad \blacksquare
\]
\end{sketchproof}

\begin{theorem}[Пряма секретність]
Компрометація довготривалих ключів $sk_C, sk_S$ після завершення
сесії не дозволяє відновити $K_s$, за умови складності CDH та
IND-CCA безпеки ML-KEM-768.
\end{theorem}

\begin{sketchproof}
Після компрометації відомо $dh_1$. Але $dh_2 = sk_S \cdot E_C$
потребує ефемерного $e_C$ (знищено); $dh_3 = e_S \cdot PK_C$~---
ефемерного $e_S$ (знищено). Аналогічно, $ss_{\mathrm{kem}}$
захищений IND-CCA ML-KEM-768: $sk_{\mathrm{kem}}$ знищено після
декапсуляції.
\[
  \mathrm{Adv}^{\mathrm{FS}} \leq
  \mathrm{Adv}_{\mathrm{CDH}} +
  \mathrm{Adv}_{\mathrm{ML\text{-}KEM}}^{\mathrm{IND\text{-}CCA}}.
  \qquad \blacksquare
\]
\end{sketchproof}

\begin{theorem}[Взаємна автентифікація]
Ймовірність підробки MAC є нехтовно малою за умови PRF-стійкості
HMAC-SHA-512:
\[
  \mathrm{Adv}^{\mathrm{auth}} \leq
  \mathrm{Adv}_{\mathrm{HMAC}}^{\mathrm{PRF}} + q_s / 2^{512}.
  \qquad \blacksquare
\]
\end{theorem}

\paragraph{Інтерпретаційне твердження про AND-властивість гібридного комбінатора.}
Нехай
\[
  \mathrm{IKM}_{\mathrm{cl}} = dh_1 \| dh_2 \| dh_3 \| dh_4,\qquad
  \mathrm{IKM}_{\mathrm{pq}} = ss_{\mathrm{kem}}.
\]
Якщо для фіксованого транскрипту $\tau$ хоча б одне з джерел
$\mathrm{IKM}_{\mathrm{cl}}$ або $\mathrm{IKM}_{\mathrm{pq}}$
залишається обчислювально нерозрізнюваним від рівномірного, тоді
\[
  \mathrm{PRK} =
  \mathrm{HKDF\text{-}Extract}(\mathrm{salt}_{\tau},
  \mathrm{IKM}_{\mathrm{cl}} \| \mathrm{IKM}_{\mathrm{pq}})
\]
може розглядатися як псевдовипадковий у межах стандартної
екстракторної інтерпретації HKDF-Extract. Операційно це підтримує таку
AND-інтерпретацію: для явного відновлення реального сеансового ключа з
транскрипту противнику потрібно
скомпрометувати \emph{обидва} джерела. Отже, йдеться не про завершене
обчислювальне доведення комбінатора, а про вмотивовану екстракторну
інтерпретацію, сумісну з прийнятою моделлю.

\begin{sketchproof}
Нехай $S$~--- те джерело серед
$\{\mathrm{IKM}_{\mathrm{cl}}, \mathrm{IKM}_{\mathrm{pq}}\}$, яке
залишається псевдовипадковим відносно зафіксованого транскрипту
$\tau$, а $X$~--- інше джерело. Тоді
$\mathrm{IKM}_{\mathrm{cl}} \| \mathrm{IKM}_{\mathrm{pq}}$ можна подати
як $S \| X$. Стандартна екстракторна властивість
HMAC-базованого HKDF-Extract~\cite{krawczyk2010} підказує, що якщо $S$
є обчислювально нерозрізнюваним від рівномірного, то й
$\mathrm{PRK}$ слід очікувати псевдовипадковим навіть за довільного
$X$. Звідси випливає груба верхня оцінка: порушення нерозрізненності
фінального ключа можливе або через злам того джерела, яке ``вижило'', або через злам
екстрактора:
\[
  \mathrm{Adv}^{\mathrm{hybrid}} \leq
  \min\!\bigl(\mathrm{Adv}_{\mathrm{CDH}}^{\mathrm{Ristretto255}},\;
  \mathrm{Adv}_{\mathrm{ML\text{-}KEM}}^{\mathrm{IND\text{-}CCA}}\bigr)
  + \mathrm{Adv}_{\mathrm{HKDF}}^{\mathrm{ext}}.
  \qquad \blacksquare
\]
\end{sketchproof}

Для геометричної інтерпретації наведеної AND-інтерпретації введемо нормалізовані
координати
\[
  x = \mathrm{Adv}_{\mathrm{CDH}}^{\mathrm{Ristretto255}},\qquad
  y = \mathrm{Adv}_{\mathrm{ML\text{-}KEM}}^{\mathrm{IND\text{-}CCA}},
\]
і розглянемо поверхню
\[
  z = \min(x, y),
\]
яка відображає головну частину верхньої межі гібридної переваги
без адитивного малого доданка
$\mathrm{Adv}_{\mathrm{HKDF}}^{\mathrm{ext}}$.

\begin{figure}[htbp]
\centering
\begin{tikzpicture}
\begin{axis}[
  width=0.92\linewidth,
  height=7.0cm,
  view={58}{28},
  domain=0:1,
  y domain=0:1,
  samples=18,
  samples y=18,
  xlabel={$x = \mathrm{Adv}_{\mathrm{CDH}}$},
  ylabel={$y = \mathrm{Adv}_{\mathrm{ML\text{-}KEM}}$},
  zlabel={$z = \min(x,y)$},
  xmin=0, xmax=1,
  ymin=0, ymax=1,
  zmin=0, zmax=1,
  xtick={0,0.25,0.5,0.75,1},
  ytick={0,0.25,0.5,0.75,1},
  ztick={0,0.25,0.5,0.75,1},
  grid=major,
  grid style={gray!20},
  tick label style={font=\small},
  label style={font=\small},
  colorbar,
  colorbar style={
    title={$z$},
    yticklabel style={font=\scriptsize},
    title style={font=\small}
  }
]
\addplot3[
  surf,
  shader=interp,
  opacity=0.95,
  colormap={ecliptix}{
    rgb255=(244,248,252)
    rgb255=(161,197,225)
    rgb255=(66,133,174)
    rgb255=(217,136,29)
  }
] {min(x,y)};
\end{axis}
\end{tikzpicture}
\caption{Просторова інтерпретація AND-моделі гібридної безпеки. По осі $x$
відкладено нормалізовану перевагу компрометації класичного шару, по осі
$y$ --- постквантового шару, а по вертикалі --- основну частину верхньої
межі $\mathrm{Adv}^{\mathrm{hybrid}} \approx \min(x,y)$. Поверхня
залишається низькою, якщо хоча б один із шарів зберігає малу перевагу
атаки, і зростає лише тоді, коли обидва джерела одночасно стають
слабкими.}
\label{fig:and-model-3d}
\end{figure}

Рисунок~\ref{fig:and-model-3d} дає інтуїтивне геометричне тлумачення: він показує,
що гібридизація не є простим ``додаванням ще одного примітива'', а
формує геометрію безпеки, у якій високий успіх противника виникає лише
в зоні одночасного ослаблення обох координатних осей.

\subsection{Стійкість до відомих атак}

\textbf{Атаки офлайн-перебирання за словником.} Захист слід читати в
межах \emph{компрометації бази даних без компрометації
\texttt{oprf\_seed}}. За цієї умови працюють три незалежні шари:
(1)~OPRF~--- без серверного ключа $k_{\mathrm{oprf}}$ противник не може
обчислити $\mathrm{rwd}$; (2)~Argon2id із параметрами MODERATE
(256~MiB, 3~ітерації)~--- на вимірювальній платформі одна спроба
коштує приблизно $T_{\mathrm{guess}} \approx 0{,}5$~с, тобто послідовний
повний перебір словника масштабується як
$C_{\mathrm{offline}} \approx |\mathcal{D}| \cdot T_{\mathrm{guess}}$;
(3)~XSalsa20-Poly1305~--- верифікація
кандидата потребує успішного автентифікованого розшифрування. Якщо ж
\texttt{oprf\_seed} скомпрометовано, це твердження не переноситься
автоматично на нову модель загроз і потребує окремого аналізу.

\textbf{Атака «людина посередині».} Транскрипт $\tau$ включає
\emph{всі} відкриті елементи, зокрема $pk_{\mathrm{kem}}$ та
$ct_{\mathrm{kem}}$. Модифікація будь-якого поля призводить до
незбіжності MAC.

\textbf{Атака повторного відтворення.} Ефемерні одноразові числа
(192~біти) та ефемерні ключі генеруються заново для кожної сесії.
Ймовірність колізії нонсів: $2^{-192}$.

\textbf{Атака «збирай зараз~--- дешифруй пізніше».} Навіть при
відновленні $dh_1,\ldots,dh_4$ алгоритмом Шора, секрет
$ss_{\mathrm{kem}}$ залишається захищеним ML-KEM-768
(Module-LWE, еквівалент AES-192).

\textbf{Атака на компрометований ключ (KCI).} Компрометація $sk_S$
не надає доступу до $dh_2 = sk_S \cdot E_C$ після закінчення сесії
(залежить від знищеного $e_C$) та до $ss_{\mathrm{kem}}$ (залежить
від знищеного $sk_{\mathrm{kem}}$).

\textbf{Часові бічні канали.} Усі порівняння MAC виконуються
рівночасовим побайтовим порівнянням без раннього виходу. Скалярне
множення на Ristretto255 та NTT в ML-KEM реалізовано в рівночасовий
спосіб.

\subsection{Формальна верифікація}

\subsubsection{Огляд верифікації}

Протокол досліджено двома незалежними символьними інструментами в
моделі Dolev--Yao: ProVerif~2.05 та Tamarin~Prover~1.10.0. Важливо,
однак, відразу зафіксувати інтерпретацію результатів. У цій роботі
символьна верифікація розглядається як \emph{доказова база} для
структурних властивостей протоколу, а не як рядок-у-рядок доведення
коректності Rust/FFI-реалізації. ProVerif використовується у двох
окремих моделях для секретності та автентифікації; Tamarin
застосовується до верифікованої сурогатної 4DH-моделі. Тому всі
підсумкові твердження мають читатися разом із межами моделі та
експериментальними матеріалами~\cite{cremers2012,basin2018}.

\begin{table}[htbp]
\caption{Підсумок формальної доказової бази}
\label{tab:verification-summary}
\centering\small
\begin{tabular}{cL{4cm}lll}
\toprule
\# & Властивість & ProVerif & Tamarin (сурогатна 4DH модель) & Rust-тести \\
\midrule
P1 & Секретність сесійного ключа & \textbf{підтверджено} & \textbf{підтверджено} (27 кр.) & 4 тести \\
P2 & Секретність пароля на дроті & \textbf{підтверджено} & \textbf{підтверджено} (3 кр.) & 6 тестів \\
P3 & Класична пряма секретність & --- & \textbf{підтверджено} (119 кр.) & 3 тести \\
P4 & Постквантова пряма секретність & --- & аргументується через P6 у межах моделі & 3 тести \\
P5a & Автентифікація ініціатора & \textbf{підтверджено} & \textbf{підтверджено} (32 кр.) & 8 тестів \\
P5b & Автентифікація респондера & \textbf{підтверджено} & \textbf{підтверджено} (17 кр.) & --- \\
P5c & Ін'єктивна взаємна автентифікація & \textbf{підтверджено} & випливає з P5a/P5b & --- \\
P6 & AND-модель гібридної безпеки & --- & \textbf{підтримано в сурогатній моделі} (29 кр.) & 4 тести \\
P7 & Стійкість до офлайн-перебору & --- & \textbf{підтримано для моделі БД без \texttt{oprf\_seed}} (4 кр.) & 6 тестів \\
--- & Коректність завершення & --- & \textbf{підтверджено} (16 кр.) & 5 тестів \\
\midrule
\textbf{Разом} & & \textbf{5/5 запитів} & \textbf{8/8 лем} & \textbf{155 пройдено, 1 пропущено} \\
\bottomrule
\end{tabular}
\end{table}

Таблиця~\ref{tab:verification-summary} навмисно поєднує символьні й
реалізаційні матеріали. Це не означає взаємозамінність моделей; радше,
йдеться про ієрархію доказовості. ProVerif краще виявляє помилки в
логіці повідомлень і кореспондентні властивості у спрощеному
обчислювальному світі необмежених сесій, тоді як Tamarin дає змогу
формалізувати станові властивості на кшталт прямої секретності та
AND-моделі. Rust-тести, своєю чергою, перевіряють, що ці властивості не
втрачаються в конкретній реалізації для широкого набору криптографічних
значень.

\begin{figure}[htbp]
\centering
\begin{tikzpicture}[
  box/.style={draw=eclblue, rounded corners=3pt, thick, fill=eclblue!8, align=center, minimum width=2.8cm, minimum height=1.05cm, font=\small},
  claim/.style={draw=eclgreen!60!black, rounded corners=3pt, thick, fill=eclgreen!10, align=center, minimum width=3.3cm, minimum height=1.15cm, font=\small},
  arr/.style={-{Latex[length=2.4mm]}, thick, draw=eclblue!80}
]
\node[box] (spec) {Протокольна\\конструкція};
\node[box, right=0.7cm of spec] (formal) {ProVerif +\\Tamarin};
\node[box, right=0.7cm of formal] (tests) {Rust-тести\\та регресії};
\node[box, below=0.9cm of spec] (wire) {Формат повідомлень\\і API};
\node[box, below=0.9cm of formal] (bench) {Criterion\\бенчмарки};
\node[box, below=0.9cm of tests] (ops) {Операційна\\межа довіри};
\node[claim, below=1.1cm of bench] (claims) {Стримані, але сильні\\наукові твердження};

\draw[arr] (spec) -- (formal);
\draw[arr] (formal) -- (tests);
\draw[arr] (wire) -- (bench);
\draw[arr] (tests) -- (ops);
\draw[arr] (formal) -- (claims);
\draw[arr] (tests) -- (claims);
\draw[arr] (bench) -- (claims);
\draw[arr] (wire) -- (claims);
\draw[arr] (ops) -- (claims);
\end{tikzpicture}
\caption{Структура доказової бази Hybrid PQ-OPAQUE: підсумкові висновки
спираються на узгодження моделі, коду, вимірювань та експлуатаційних
меж, а не на один ізольований матеріал.}
\label{fig:evidence-stack}
\end{figure}

\subsubsection{ProVerif 2.05}

Інструмент: ProVerif~2.05 (OCaml~5.4.0), macOS Darwin~25.2.0. Використано
дві моделі: \texttt{formal/hybrid\_pq\_opaque.pv} для секретності та
\texttt{formal/hybrid\_pq\_opaque\_auth.pv} для автентифікації.
Поділ на дві моделі є не технічною дрібницею, а методологічним
рішенням: одна й та сама насичена модель не давала стабільної збіжності
для всіх запитів одночасно. Тому результати ProVerif потрібно читати як
сукупність двох сумісних, але різних поглядів на протокол.

Найбільш інформативним для властивості взаємної автентифікації є запит
P5c (ін'єктивна взаємна відповідність):

\begin{lstlisting}
query pkC: point, pkS: point, sk: key;
  inj-event(ServerCompletesAuth(pkS, pkC, sk))
  ==> inj-event(ClientCompletesAuth(pkC, pkS, sk)).
RESULT: true
\end{lstlisting}

Логічний результат \texttt{true} означає, що завершення сесії на стороні
респондера не може бути пояснене повторним програванням одного й того ж
успішного трасування для кількох клієнтських завершень. Водночас цей
результат не слід тлумачити як обчислювальне доведення всієї реалізації:
в моделі відсутні бічні канали, варіативність формату помилок,
витоки, пов'язані з часовими характеристиками, та низка мережевих деталей, які залишаються предметом
реалізаційних тестів.

\subsubsection{Tamarin Prover 1.10.0}

Інструмент: Tamarin~1.10.0 (Maude~3.4). Верифікована модель:
\texttt{formal/hybrid\_pq\_opaque\_verified.spthy}. Архівований лог
\texttt{formal/logs/tamarin\_full\_proof.log} фіксує загальний час
обробки \textbf{22,83~секунди} для всіх восьми лем.

\begin{table}[htbp]
\caption{Підтверджені леми Tamarin для сурогатної 4DH-моделі}
\label{tab:tamarin-lemmas}
\centering\small
\begin{tabularx}{\linewidth}{L{4.5cm}rX}
\toprule
Лема & Кроки & Ключова властивість \\
\midrule
\texttt{session\_key\_secrecy} (P1) & 27 & Без подій корупції противник не деривує сесійний ключ \\
\texttt{password\_secrecy} (P2) & 3 & \texttt{\textasciitilde{}pwd} не потрапляє до \texttt{Out()} \\
\texttt{forward\_secrecy} (P3) & 119 & Постсесійна компрометація LTK не розкриває ефемерні секрети \\
\texttt{auth\_initiator} (P5a) & 32 & Успішний KE3 вимагає попереднього \texttt{RunR} \\
\texttt{auth\_responder} (P5b) & 17 & Успішний \texttt{Finish} вимагає валідного KE3 \\
\texttt{and\_model} (P6) & 29 & Без \texttt{RevEph} неможливо відновити 4DH та KEM ss разом \\
\texttt{dictionary\_resistance} (P7) & 4 & Компр. БД розкриває запис, але не OPRF-ключ \\
\texttt{completion} (коректність) & 16 & Існує трейс з успішним завершенням \\
\bottomrule
\end{tabularx}
\end{table}

Суттєво, що верифікована Tamarin-модель не використовує вбудовану
рівняльну теорію \texttt{diffie-hellman} для повної 4DH-конструкції.
Натомість вона замінює 4DH на абстрактний секрет
\texttt{h(<"4dh", skC, skS, ekC, ekS>)} й перевіряє, які події
компрометації потрібні противнику для виведення сесійного ключа. Це
істотно полегшує автоматизоване доведення, але саме
тому її результати слід трактувати як сурогатну символьну аргументацію,
а не як формальний доказ конкретної реалізації кожної операції
Ristretto255.

Центральною для подальшої інтерпретації є лема P6 (AND-модель):
\begin{lstlisting}
All C S sk #i. SKC(C, S, sk) @ #i
  & not(Ex #j. RevEph(C) @ #j)
  & not(Ex #j. RevEph(S) @ #j)
  ==> not(Ex #j. K(sk) @ #j)
Result: verified (29 steps)
\end{lstlisting}

Лема допускає компрометацію LTK, але забороняє розкриття ефемерних
ключів. Навіть знаючи $sk_A$ та $sk_R$, противник не може обчислити
\texttt{combine(dh\_secret, kem\_ss)} без ефемерних \texttt{\textasciitilde{}ek} та
\texttt{\textasciitilde{}ksk}. У межах цієї сурогатної моделі це
підтримує інтерпретацію AND-моделі: для виведення сесійного ключа
потрібне одночасне ослаблення \emph{обох} шарів. Саме ця лема дає
лише опосередковану аргументацію на користь P4; вона не замінює окремої
формалізації постквантової прямої секретності.

\subsubsection{Межі інтерпретації результатів}

Щоб уникнути надмірної інтерпретації формальних результатів, зафіксуємо
чотири ключові обмеження.
\begin{enumerate}
  \item \textbf{Сурогатне подання 4DH.} Верифікована модель Tamarin абстрагує
        4DH-компонент до однієї функції знання. Такий крок полегшує
        доведення, але не замінює аналіз конкретних алгебраїчних властивостей
        групових операцій.
  \item \textbf{Розділені моделі ProVerif.} Секретність і автентифікація
        досліджуються не в одній універсальній моделі, а в двох
        окремих матеріалах. Отже, результат 5/5 слід трактувати як
        сукупність двох узгоджених моделей, а не як монолітне доведення.
  \item \textbf{Межа застосовності тверджень про офлайн-стійкість.} Лема P7 і
        відповідні Rust-тести покривають сценарій компрометації бази
        даних \emph{без} компрометації \texttt{oprf\_seed}. Якщо
        \texttt{oprf\_seed} розкрито, властивість офлайн-стійкості
        потребує окремого аналізу і не повинна автоматично переноситися
        з моделі.
  \item \textbf{Розрив між моделлю та реалізацією.} Символьна верифікація не охоплює
        бічні канали, поведінку компілятора, властивості виділення
        пам'яті, API-помилки FFI чи вразливості зовнішніх крейтів. Саме
        тому результати формальної верифікації в роботі доповнюються
        тестами й еталонними вимірюваннями, а не підміняють їх.
\end{enumerate}

% ============================================================
\section{Реалізація та експериментальна оцінка}
% ============================================================

\subsection{Архітектура Rust-реалізації}

Протокол реалізовано мовою \textbf{Rust~1.93.1} у вигляді
багатокрейтного простору проєктів. Вибір Rust обумовлений:
гарантіями безпеки пам'яті на рівні системи типів (відсутність
нульових покажчиків, відсутність використання після звільнення,
інспектор запозичень); підтримкою \texttt{\#![forbid(unsafe\_code)]}
у крейтах ядра; трейтом \texttt{zeroize} для безпечного затирання
ключів у пам'яті.

\begin{table}[htbp]
\caption{Метрики вихідного коду простору проєктів}
\label{tab:code-metrics}
\centering\small
\begin{tabularx}{\linewidth}{L{2.5cm}L{4cm}rL{3.5cm}}
\toprule
Крейт & Призначення & РВК & Залежності \\
\midrule
\texttt{opaque-core} & Криптографічні примітиви & 1509 & curve25519-dalek, ml-kem, argon2 \\
\texttt{opaque-agent} & Ініціатор (клієнт) & 899 & opaque-core \\
\texttt{opaque-relay} & Респондер (сервер) & 907 & opaque-core \\
\texttt{opaque-ffi} & Прошарок FFI & 2020 & opaque-agent, opaque-relay \\
\midrule
\textbf{Разом (основний код)} & & \textbf{5335} & \\
Тести & & 3961 & proptest, власні регресійні сценарії \\
Еталонні вимірювання & & 1460 & Criterion \\
Формальні моделі & & 1806 & Tamarin, ProVerif \\
\bottomrule
\end{tabularx}
\end{table}

\subsection{Безпека пам'яті та затирання секретів}

Для криптографічної системи важливо не лише те, які примітиви
використовуються, а й те, як вони живуть у пам'яті процесу. У поточному
просторі проєктів крейти \texttt{opaque-core}, \texttt{opaque-agent} та
\texttt{opaque-relay} компілюються з директивою
\texttt{\#![forbid(unsafe\_code)]}; небезпечний код локалізовано лише в
FFI-межі, де взаємодія із сирими покажчиками неминуча. Така
архітектура зменшує ризик помилок класу use-after-free, подвійного
звільнення або порушення інваріантів довжини буфера.

Окремий акцент зроблено на життєвому циклі секретів. Реалізація широко
використовує \texttt{Zeroize}, \texttt{ZeroizeOnDrop} та обгортку
\texttt{Zeroizing}, завдяки чому проміжні значення не залишаються в
пам'яті довше, ніж цього вимагає протокол. Зокрема, затираються:
\begin{itemize}
  \item пароль користувача та OPRF-осліплюючий скаляр після фіналізації
        OPRF;
  \item ефемерний секрет ML-KEM після декапсуляції;
  \item класичний 4DH-матеріал, transcript hash і PRK після деривації
        похідних ключів;
  \item контекстний хеш \texttt{account\_id} після завершення KE2/KE3;
  \item станові структури ініціатора й респондера при протермінуванні
        або будь-якому завершенні з безпечним відхиленням.
\end{itemize}
Це істотно для зменшення ризику посткомпрометаційного
витоку через дампи пам'яті, аварійні логи або повторне використання
буферів.

\subsection{Тестова база та методологія}

Набір тестів будується за принципом відповідності тестових модулів
властивостям безпеки, встановленим у Розділі~3. Тестування
поділяється на три методологічні класи:
\begin{itemize}
  \item \emph{Детерміністичні тести}~--- конкретні протокольні
        сценарії з фіксованими або псевдовипадковими вхідними даними;
  \item \emph{Інтеграційні тести}~--- повний наскрізний цикл
        реєстрація~$\to$ автентифікація;
  \item \emph{Тести на основі властивостей} (бібліотека
        \texttt{proptest})~--- довільні вхідні дані за визначеними
        стратегіями.
\end{itemize}

\begin{table}[htbp]
\caption{Актуальна структура тестового набору}
\label{tab:test-distribution}
\centering\small
\begin{tabularx}{\linewidth}{L{4.5cm}L{4cm}rL{2.5cm}}
\toprule
Модуль & Властивість & Кількість & Клас \\
\midrule
\texttt{opaque-agent} модульні тести & Життєвий цикл стану ініціатора & 3 & Модульний \\
\texttt{integration} & Наскрізний цикл і сумісність мережевого формату & 9 & Інтеграційний \\
\texttt{security\_properties} & Деталізовані класи безпеки P1--P7 & 34 & Детерміністичний \\
\texttt{security\_proptest} & P1, P2, P5 на довільних даних & 5 & Властивісний \\
\texttt{security\_regressions} & Відомі регресії на межі ініціатора та релейного вузла & 8 & Регресійний \\
\texttt{crypto\_tests + envelope\_tests} & Базові криптографічні примітиви & 38 & Модульний \\
\texttt{oprf\_tests + pq\_kem\_tests} & OPRF і KEM-коректність & 21 & Модульний \\
\texttt{protocol\_tests} & Серіалізація та версійність & 17 & Модульний \\
\texttt{opaque-ffi} тести & FFI-контракт і помилки API & 4 & Інтеграційний \\
\texttt{opaque-relay} тести & Серверний рівень і стан респондера & 17 & Модульний \\
\midrule
\textbf{Разом} & & \textbf{156 визначено} & \textbf{155 пройдено, 1 пропущено} \\
\bottomrule
\end{tabularx}
\end{table}

\begin{lstlisting}
cargo test --workspace
test result: ok. 155 passed; 0 failed; 1 ignored
(Rust 1.93.1, debug profile; ignored test = slow compromise-boundary regression)
\end{lstlisting}

\subsection{Трасування безпекових вимог до тестових матеріалів}

Для прикладної криптографічної реалізації важливо не лише мати великий
набір тестів, а й показати, \emph{які саме} властивості протоколу
перевіряє кожен тип матеріалів. Інакше тестова база легко
перетворюється на нерозрізнений набір сценаріїв без явного зв'язку з
науковими твердженнями роботи. Таблиця~\ref{tab:traceability} фіксує
відповідність між вимогами Розділу~3, формальними моделями та
реалізаційними тестами.

\begin{table}[htbp]
\caption{Трасування вимог R1--R5 до формальних моделей і тестових матеріалів}
\label{tab:traceability}
\centering\small
\begin{tabularx}{\linewidth}{L{0.9cm}L{3.4cm}L{4.2cm}X}
\toprule
Вимога & Формальна опора & Реалізаційні матеріали & Інтерпретація \\
\midrule
R1 & Твердження про межі офлайн-перебирання, пряму секретність і взаємну автентифікацію; ProVerif (секретність, автентифікація) & \texttt{security\_properties}, \texttt{security\_proptest}, \texttt{integration}, \texttt{envelope\_tests} & Перевіряються парольна таємність, завершення сесії, цілісність MAC і коректність роботи конверта \\
R2 & Інтерпретаційне твердження про AND-властивість; Tamarin P6 & \texttt{pq\_kem\_tests}, \texttt{security\_regressions}, \texttt{protocol\_tests} & Перевіряється включення KEM-компонента до транскрипту й відсутність мовчазного класичного режиму \\
R3 & Інтерпретаційне твердження про AND-властивість; Tamarin \texttt{and\_model} & \texttt{security\_properties}, тести \texttt{opaque-relay} & Властивість AND-моделі підкріплюється як символьно, так і перевіркою життєвого циклу серверного стану \\
R4 & Формальна структура протоколу KE1--KE2--KE3 & \texttt{protocol\_tests}, \texttt{integration}, еталонні вимірювання Criterion & Підтверджуються незмінні 1,5~RTT, версіонування формату повідомлень та кількісні накладні витрати \\
R5 & Редукційні ескізи та вимоги безпечного відхилення & тести \texttt{opaque-ffi}, модульні тести \texttt{opaque-agent}, бенчмарки, перевірки безпечного затирання & Перевіряються API-контракт, коректне завершення на помилках, безпечне очищення стану та відтворюваність продуктивності \\
\bottomrule
\end{tabularx}
\end{table}

Така схема має два методологічні наслідки. По-перше, вона зменшує
ризик підміни понять, коли позитивний результат мікротесту помилково
трактується як підтвердження складної протокольної властивості.
По-друге, вона показує, що жоден окремий тип матеріалів не є
достатнім сам по собі: символьні моделі покривають абстрактну логіку,
тоді як Rust-тести виявляють помилки серіалізації, життєвого циклу
стану, перевірки меж вхідних даних і FFI-контракту, які принципово відсутні у
формальних моделях.

\subsection{Методологія еталонних вимірювань}

\textbf{Апаратна платформа:} Apple~M1~Pro (8~ВП + 2~ЕЕ ядра,
16~ГБ LPDDR5, 3,228~ГГц), macOS~26.2, Rust~1.93.1.

\textbf{Фреймворк:} Criterion~0.5.1~\cite{heyer2023} застосовує
метод бутстрапу Welch: $N_s = 100$ незалежних вибірок, кожна з яких
є вибірковим середнім $\overline{t}_i$ по $n_i$~ітерацій; ітерації
підбираються автоматично до відносної похибки $\epsilon < 0{,}01$.
Вибіркове середнє $\hat{\mu}$ та 95\%-й довірчий інтервал (ДІ)
отримуються методом бутстрапу з 100\,000~перевибірок. Фаза
прогріву~--- 3~с.

\textbf{Умови:} \texttt{--release}, \texttt{opt-level~=~3},
\texttt{lto~=~false}; детермінований однопотоковий запуск.

\subsection{Процедура відтворення результатів}

Окремою вимогою до сучасної прикладної криптографічної статті є
можливість відтворити не лише окремі цифри, а й сам процес отримання
доказової бази. У межах цієї роботи відтворення поділяється на три
рівні: (1) запуск тестів реалізації; (2) запуск еталонних вимірювань;
(3) перегляд або повторний запуск формальних моделей. Такий поділ
важливий тому, що час виконання й інструментальні залежності для цих
рівнів істотно відрізняються.

\begin{table}[htbp]
\caption{Мінімальна процедура відтворення дослідницьких матеріалів}
\label{tab:artifact-repro}
\centering\small
\begin{tabularx}{\linewidth}{L{2.6cm}L{4.0cm}L{2.4cm}X}
\toprule
Шар & Команда / матеріал & Типовий час & Очікуваний результат \\
\midrule
Тести реалізації & \texttt{cargo test --workspace} & Хвилини & 155~пройдених тестів, 1~пропущений сценарій граничної регресії \\
Локальні регресійні тести безпеки & \texttt{cargo test -p opaque-agent --test security\_regressions} & Секунди & Перевірка безпечного відхилення на межі між модулями agent і relay \\
Еталонні вимірювання & \texttt{cargo bench --workspace} & Від хвилин до десятків хвилин & Звіти Criterion для мікро-, протокольних сценаріїв і сценаріїв пропускної здатності \\
ProVerif & \texttt{proverif formal/hybrid\_pq\_opaque.pv} та \texttt{proverif formal/hybrid\_pq\_opaque\_auth.pv} & Секунди/хвилини & Підтвердження секретності та відповідності в окремих моделях \\
Tamarin & Архівовані логи в \texttt{formal/logs/}; за потреби --- \texttt{./formal/run-tamarin.sh} & Від швидкого перегляду до тривалого обчислення & Відтворення або аудит сурогатної 4DH-моделі та меж її інтерпретації \\
\bottomrule
\end{tabularx}
\end{table}

Принципово, що стаття не вимагає від рецензента одразу повторювати весь
обсяг формальної роботи з нуля. Базовий аудит може починатися з
архівованих логів, структури тестів і відтворення бенчмарків. Повторний
повний запуск Tamarin є цінним, але не повинен бути єдиною точкою
довіри. Саме така багатошарова організація дослідницьких матеріалів робить дослідження
зручним для зовнішньої перевірки: різні члени журі або комісії можуть
зайти в нього з того рівня технічної глибини, який для них доступний.

\subsection{Еталонні вимірювання криптографічних примітивів}

\begin{table}[htbp]
\caption{Еталонні вимірювання криптографічних примітивів}
\label{tab:micro-bench}
\centering\small
\begin{tabular}{llrr}
\toprule
Примітив & Операція & $\hat{\mu}$ & 95\% ДІ \\
\midrule
Ristretto255 & Генерація ключової пари & 16,74~мкс & [16,03; 17,84] \\
Ristretto255 & Одиничне DH & 46,41~мкс & [45,94; 46,98] \\
Ristretto255 & 3DH (еталон) & 119,04~мкс & [117,10; 122,14] \\
Ristretto255 & \textbf{4DH (протокол)} & $\approx\mathbf{164}$~\textbf{мкс} & [$\approx$162; $\approx$167] \\
ML-KEM-768 & Генерація ключової пари & 41,05~мкс & [40,69; 41,68] \\
ML-KEM-768 & Інкапсуляція & 33,62~мкс & [32,89; 34,95] \\
ML-KEM-768 & Декапсуляція & 38,55~мкс & [37,65; 40,11] \\
ML-KEM-768 & \textbf{Повний раунд} & \textbf{122,01~мкс} & [119,32; 126,96] \\
OPRF & Засліплення & 56,16~мкс & [54,37; 59,36] \\
OPRF & Обчислення сервера & 47,59~мкс & [44,97; 52,45] \\
OPRF & Фіналізація & 65,68~мкс & [64,17; 67,89] \\
HKDF & Витяг & 1,32~мкс & [1,312; 1,319] \\
HKDF & Розширення (64~байти) & 1,21~мкс & [1,181; 1,245] \\
HMAC-SHA-512 & 256~байт & 1,61~мкс & [1,596; 1,641] \\
XSalsa20-Poly1305 & Шифрування 96~байт & 440,79~нс & [430,20; 459,46] \\
XSalsa20-Poly1305 & Розшифрування 96~байт & 613,42~нс & [594,15; 636,12] \\
Argon2id & \textbf{MODERATE} & \textbf{510,86~мс} & [490,31; 533,80] \\
Гібридний комбінатор & HKDF-Extract (160~байт) & 1,48~мкс & [1,453; 1,508] \\
\bottomrule
\end{tabular}
\end{table}

\textbf{Аналітичні спостереження.}
\begin{enumerate}
  \item \emph{Порівнянна вартість класичного та постквантового шарів.} Повний
        раунд ML-KEM-768 ($\hat{\mu} = 122$~мкс) та 4DH
        ($\hat{\mu} \approx 164$~мкс) мають близькі порядки
        величин. Відношення постквантових до класичних накладних витрат~---
        лише $\times 1{,}06$, що не змінює асимптотичний клас
        складності протоколу.
  \item \emph{Абсолютне домінування KSF.} Argon2id
        ($\hat{\mu} = 510{,}86$~мс) перевищує повний раунд
        ML-KEM-768 у $4\,175$~разів. Це свідоме проєктне рішення.
  \item \emph{Практично незначна вартість гібридного комбінатора.}
        \texttt{combine\_key\_material} ($\hat{\mu} = 1{,}48$~мкс)
        вносить менше 0,001\% у загальний наскрізний час.
\end{enumerate}

\subsection{Еталонні вимірювання протокольного рівня}

\begin{table}[htbp]
\caption{Еталонні вимірювання протокольних функцій}
\label{tab:protocol-bench}
\centering\small
\begin{tabularx}{\linewidth}{L{2cm}L{3cm}rrX}
\toprule
Фаза & Функція & $\hat{\mu}$ & 95\% ДІ & Домінуючий компонент \\
\midrule
Реєстрація & \texttt{create\_request} & 72,46~мкс & [72,15; 72,77] & Засліплення OPRF \\
Реєстрація & \texttt{create\_response} & 45,92~мкс & [45,72; 46,16] & Обчислення OPRF \\
Реєстрація & \texttt{finalize} & \textbf{479,58~мс} & [472,11; 487,75] & \textbf{Argon2id} \\
Автентифікація & \texttt{generate\_ke1} & 118,83~мкс & [116,59; 121,76] & OPRF~+ генерація ключів ML-KEM \\
Автентифікація & \texttt{generate\_ke2} & 297,03~мкс & [294,81; 299,69] & 4DH~+ інкапсуляція ML-KEM \\
Автентифікація & \texttt{generate\_ke3} & \textbf{492,34~мс} & [478,95; 509,96] & \textbf{Argon2id}~+ декапсуляція ML-KEM \\
Автентифікація & \texttt{responder\_finish} & 1,24~мкс & [1,174; 1,323] & Перевірка MAC \\
\midrule
Повний & \texttt{authentication\_e2e} & \textbf{481,43~мс} & [471,73; 495,25] & \textbf{Argon2id} \\
\bottomrule
\end{tabularx}
\end{table}

Числа в Табл.~\ref{tab:protocol-bench} не є адитивними. Значення
\texttt{generate\_ke3} та \texttt{authentication\_e2e} отримано в
окремих сценаріях Criterion із різним контуром підготовки стану та
різною структурою ітерації. Саме тому середнє для повного сценарію може
виявитися трохи нижчим за середнє для ізольованої фази; наведені
95\% довірчі інтервали при цьому суттєво перекриваються. У цій роботі такі
вимірювання використовуються насамперед для порівняння порядків
величин і домінуючих компонент, а не для пофазового арифметичного
складання.

\textbf{Аналіз \texttt{generate\_ke2} = 297~мкс.} Час виконання:
\begin{equation}
  4\mathrm{DH} \approx 164 + \mathrm{encaps} \approx 34 +
  \mathrm{HKDF\text{-}Expand} \approx 6 + \mathrm{MAC} \approx 2 +
  \delta_{\mathrm{op}} \approx 91~\text{мкс}.
\end{equation}
Четвертий добуток $dh_4 = e_S \cdot E_C$ додає явне зв'язування обох
ефемерних компонент і пояснює зростання часу порівняно з 3DH приблизно
на 41~мкс.

\subsection{Оцінка пропускної здатності респондера}

\begin{table}[htbp]
\caption{Оцінка пропускної здатності вузла-респондера}
\label{tab:throughput}
\centering\small
\begin{tabular}{lrl}
\toprule
Сценарій & $\hat{\mu}$ & Пропускна здатність \\
\midrule
\texttt{finish\_only}\footnote{Сценарій \texttt{finish\_only} вимірює лише фінальну фазу завершення: верифікацію MAC та витяг сесійного ключа через \texttt{responder\_finish}. Фаза генерації KE2 (4DH + інкапсуляція ML-KEM-768 + деривація ключів HKDF) виконується на етапі підготовки еталонного вимірювання та оцінюється окремо сценарієм \texttt{ke2\_only}.} (без Argon2id) & 760,82~нс & \textbf{1,31~млн викликів/с} \\
\texttt{ke2\_only} (без Argon2id) & 254,26~мкс & 3,93~тис. автентиф./с \\
\bottomrule
\end{tabular}
\end{table}

Таблицю~\ref{tab:throughput} потрібно інтерпретувати обережно. Значення
1,31~млн/с характеризує \emph{лише} виклик \texttt{responder\_finish}
після того, як KE2 вже згенеровано і стан респондера підготовлено.
Повна серверна вартість однієї автентифікації визначається насамперед
фазою \texttt{ke2\_only}. Якщо скласти обидва етапи, одержуємо
\[
  \frac{1}{254{,}26~\mu s + 0{,}76082~\mu s} \approx 3{,}92 \times 10^3
  \text{ автентифікацій/с на ядро}.
\]
Отже, коректний висновок полягає не в «мільйоні автентифікацій на
секунду», а в тому, що \texttt{responder\_finish} є практично
безвартісним порівняно з генерацією KE2 і що серверна пропускна
здатність обмежується саме криптографічною фазою KE2, а не завершальним
MAC-підтвердженням.

\subsection{Кількісна оцінка накладних витрат постквантового компонента}

\begin{table}[htbp]
\caption{Накладні витрати ML-KEM-768 за фазами AKE (без Argon2id)}
\label{tab:pq-overhead}
\centering\small
\resizebox{\linewidth}{!}{%
\begin{tabular}{lrrrrl}
\toprule
Фаза AKE & Класична (4DH) & Гібридна (4DH+KEM) & Абс. надлишок & Відн. надлишок & Постквантова операція \\
\midrule
KE1 (ініціатор) & 73,2~мкс & 114,5~мкс & \textbf{$+41{,}3$~мкс} & $+56{,}4\%$ & Генерація ключів KEM \\
KE2 (респондер, без KSF) & 219,0~мкс & 239,5~мкс & \textbf{$+20{,}4$~мкс} & $+9{,}3\%$ & Інкапсуляція~+ комбінатор \\
KE3 (ініціатор, без KSF) & 234,0~мкс & 267,6~мкс & \textbf{$+33{,}2$~мкс} & $+14{,}2\%$ & Декапсуляція~+ комбінатор \\
\midrule
\textbf{Повний AKE (без KSF)} & \textbf{528,2~мкс} & \textbf{687,5~мкс} & \textbf{$+159{,}3$~мкс} & \textbf{$+30{,}1\%$} & ГКП + Інк.~+ Дек. \\
\bottomrule
\end{tabular}%
}
\end{table}

\begin{figure}[htbp]
\centering
\begin{tikzpicture}
\begin{axis}[
  width=0.95\linewidth,
  height=6.2cm,
  ybar,
  bar width=9pt,
  ymin=0,
  ymax=300,
  ylabel={Мікросекунди},
  symbolic x coords={KE1,KE2,KE3},
  xtick=data,
  enlarge x limits=0.24,
  ymajorgrids=true,
  grid style={gray!20},
  legend style={at={(0.5,1.02)}, anchor=south, legend columns=2, draw=none, font=\small},
  tick label style={font=\small},
  label style={font=\small}
]
\addplot[fill=eclblue!75, draw=eclblue] coordinates {
  (KE1,73.2) (KE2,219.0) (KE3,234.0)
};
\addplot[fill=eclred!70, draw=eclred!80!black] coordinates {
  (KE1,114.5) (KE2,239.5) (KE3,267.6)
};
\legend{Класична AKE,Гібридна AKE}
\end{axis}
\end{tikzpicture}
\caption{Порівняння класичного та гібридного AKE без Argon2id:
додавання ML-KEM-768 збільшує час кожної фази, але не змінює її
порядок величини і не робить постквантовий компонент домінуючим.}
\label{fig:pq-overhead-phases}
\end{figure}

Рисунок~\ref{fig:pq-overhead-phases} добре показує, чому
гібридизація є практично прийнятною: навіть після додавання
ML-KEM-768 жодна фаза без KSF не переходить в інший порядок величини за латентністю, а приріст
залишається в межах десятків мікросекунд.

\textbf{Розкладання постквантових накладних витрат:}
\begin{equation}
  \Delta_{PQ} \approx
  \underbrace{41{,}4}_{\text{ген. кл.}} +
  \underbrace{33{,}8}_{\text{інкапс.}} +
  \underbrace{39{,}2}_{\text{декапс.}} +
  \underbrace{2 \times 1{,}48}_{\text{комб.}} +
  \underbrace{\delta_{\mathrm{алок.}}}_{\approx 42}
  \approx 119{,}3 + 42 = 161~\text{мкс},
\end{equation}
де $\delta_{\mathrm{алок.}}$~--- витрати алокації буферів
ML-KEM-768: $|pk_{\mathrm{kem}}| \approx 1{,}2$~КБ,
$|sk_{\mathrm{kem}}| \approx 2{,}4$~КБ,
$|ct_{\mathrm{kem}}| \approx 1{,}1$~КБ. Порівняно з Argon2id
(510,86~мс), постквантовий надлишок становить близько $0{,}03\%$ типового
наскрізного часу.

\subsection{Порівняльний аналіз гібридних протоколів}

\begin{table}[htbp]
\caption{Порівняння гібридних постквантових протоколів у зіставних для цієї роботи аспектах}
\label{tab:protocol-comparison}
\centering\small
\resizebox{\linewidth}{!}{%
\begin{tabular}{lllll}
\toprule
Властивість & TLS~1.3 Hybrid & Signal PQXDH & WireGuard~PQ & \textbf{Hybrid PQ-OPAQUE} \\
\midrule
Класична основа & ECDH (X25519) & X3DH (X25519) & Noise IK (X25519) & \textbf{4DH (Ristretto255)} \\
Постквантовий компонент & ML-KEM-768 & ML-KEM-1024 & ML-KEM-768 & \textbf{ML-KEM-768} \\
Комбінування ключового матеріалу & Гібридне поєднання & Гібридне поєднання & Гібридне поєднання & \textbf{Гібридне поєднання з AND-інтерпретацією} \\
Тип автентифікації & PKI-серт. & Пакет поперед. кл. & Статичні ключі & \textbf{PAKE (OPRF~+~пароль)} \\
Сценарій пасивної компрометації БД & Не застосовується & Не застосовується & Не застосовується & \textbf{Підтримується у моделі БД} \\
Захист від офл. перебирання & Не застосовується & Не застосовується & Не застосовується & \textbf{Підтримується у межах моделі} \\
Пряма секретність (класична) & Так & Так & Ні & \textbf{Так} \\
Постквантова пряма секретність & Так & Так & Ні & \textbf{Підтримується у межах моделі} \\
Формальна доказова база & Окремі результати & Окремі результати & Не заявлено & \textbf{ProVerif~+ Tamarin (симв.)} \\
Орієнтовний обсяг рукостискання & $\approx$3200~байт & $\approx$2400~байт & $\approx$2800~байт & \textbf{2715~байт у мережевому форматі} \\
Кількість раундів & 1~RTT & 0~RTT & 1~RTT & \textbf{1,5~RTT} \\
\bottomrule
\end{tabular}%
}
\end{table}

Наведене порівняння має кілька важливих інтерпретаційних меж. По-перше,
числа для Hybrid PQ-OPAQUE подано у мережевому форматі з 1-байтовими
префіксами версії, що відрізняє їх від поширеної практики рахувати
лише корисні дані. Для TLS~1.3 Hybrid, Signal PQXDH і WireGuard~PQ
наведено літературні орієнтовні оцінки порядку величини, а не власний
уніфікований перерахунок байтового обсягу з однаковими правилами для
транспортного кадрування, версійних префіксів і службових полів. Отже,
строге байтове зіставлення коректне лише для нашого протоколу, тоді як
зовнішні числа слід читати як порівняння порядків величин.
По-друге, графа «Формальна доказова база» не означає
наявності обчислювального доведення всієї реалізації; для нашого
протоколу йдеться про символьну доказову базу, а не про повністю
машинно перевірене доведення коду. По-третє, рядки
«Комбінування ключового матеріалу», «Сценарій пасивної компрометації БД»,
«Захист від офл. перебирання» та «Постквантова пряма секретність» є стислим порівняльним резюме, а не
повними обчислювальними твердженнями. Для Hybrid PQ-OPAQUE їх слід
читати лише в межах сценарію компрометації бази даних \emph{без}
компрометації \texttt{oprf\_seed}; крім того, AND-інтерпретація
означає сумісність із такою інтерпретацією на рівні комбінатора
ключового матеріалу, а не завершене доведення комбінатора в усіх моделях.

\subsection{Формальні моделі як відтворювані дослідницькі матеріали}

Важливою частиною дослідження є не лише опис властивостей, а й
публікація самих формальних моделей як самостійних дослідницьких матеріалів.
Таблиця~\ref{tab:formal-models} систематизує модельний пакет, що
супроводжує реалізацію.

\begin{table}[htbp]
\caption{Формальні моделі, включені до пакета дослідницьких матеріалів}
\label{tab:formal-models}
\centering\small
\begin{tabularx}{\linewidth}{L{4.6cm}L{2.1cm}rX}
\toprule
Файл & Інструмент & РВК & Призначення \\
\midrule
\texttt{hybrid\_pq\_opaque.spthy} & Tamarin & 533 & Повна 4DH-модель, орієнтир для аналізу, не сходиться практично \\
\texttt{hybrid\_pq\_opaque\_ascii.spthy} & Tamarin & 533 & ASCII-подання повної моделі для зовнішнього перегляду \\
\texttt{hybrid\_pq\_opaque\_verified.spthy} & Tamarin & 193 & Верифікована сурогатна модель 4DH \\
\texttt{hybrid\_pq\_opaque.pv} & ProVerif & 296 & Секретність ключа та пароля \\
\texttt{hybrid\_pq\_opaque\_auth.pv} & ProVerif & 251 & Автентифікація та ін'єктивна відповідність \\
\midrule
\textbf{Разом} & & \textbf{1806} & \\
\bottomrule
\end{tabularx}
\end{table}

Таке розділення дослідницьких матеріалів має два наслідки. По-перше, воно робить
статтю відтворюваною: сторонній дослідник може побачити не лише
кінцеву таблицю результатів, а й конкретні моделі, які привели до
цих результатів. По-друге, воно дисциплінує формулювання тверджень:
якщо одна властивість доводиться в окремій моделі, а інша --- в
сурогатній, текст статті має прямо про це повідомляти. Саме тому в цій роботі
розділено поняття ``символьна верифікація моделей'' та ``формальна доказова
база''.

\subsection{Перевірка вхідних даних і раннє відхилення некоректних станів}

Оскільки протокол орієнтовано на прикладне використання через Rust API і
FFI, окрему роль відіграє рання валідація вхідних даних. Реалізація
відхиляє некоректні або потенційно небезпечні значення до початку
основних криптографічних обчислень. До цієї категорії належать:
\begin{itemize}
  \item неправильні довжини KE1/KE2/KE3 та реєстраційних повідомлень;
  \item невідомі префікси версії мережевого формату;
  \item тотожні або неканонічні точки Ristretto255;
  \item порожній \texttt{account\_id} та інші некоректні параметри API;
  \item повторне використання або протермінування об'єктів стану.
\end{itemize}

З наукової точки зору ці перевірки важливі тому, що вони зменшують
розрив між символьною моделлю та реалізацією. У формальній моделі
противник часто подає довільні терми; у реалізації цю свободу треба
відтворити безпечно через перевірку мережевого формату, валідності точок,
життєвого циклу стану та обмежень API. Саме тому політика раннього
відхилення не є ``інженерною дрібницею'', а частиною реалізаційної аргументації
безпеки.

% ============================================================
\section{Обговорення результатів}
% ============================================================

\subsection{Відповідність встановленим вимогам}

\begin{table}[htbp]
\caption{Оцінка відповідності протоколу встановленим вимогам у межах прийнятої моделі}
\label{tab:requirements}
\centering\small
\begin{tabularx}{\linewidth}{L{0.8cm}L{2.2cm}L{2.9cm}X}
\toprule
Вимога & Опис & Статус & Обґрунтування \\
\midrule
R1 & OPAQUE-властивості & Підтримано & aPAKE-цілі модифікованої схеми, роль OPRF, Argon2id і конверта збережено; 155 пройдених тестів \\
R2 & Постквантова стійкість & Підтримано & ML-KEM-768 IND-CCA2 (FIPS~203); постквантовий компонент включено до транскрипту та схеми виведення ключів \\
R3 & AND-інтерпретація & Підтримано інтерпретаційно & $\mathrm{HKDF}(dh_{1..4} \| ss_{\mathrm{kem}})$; інтерпретаційне твердження та Tamarin P6 підтримують AND-інтерпретацію в прийнятій моделі \\
R4 & Мінімальні накладні витрати & Підтверджено вимірюваннями & $+159{,}3$~мкс процесорного часу (без KSF); $+2272$~байти у мережевому форматі; 0~додаткових раундів; близько $0{,}03\%$ \\
R5 & Реалізованість & Підтверджено дослідницькими матеріалами & 5335~РВК Rust, 4~крейти, 155~пройдених тестів, Criterion; відтворюваний пакет дослідницьких матеріалів \\
\bottomrule
\end{tabularx}
\end{table}

У Табл.~\ref{tab:requirements} статуси означають ступінь підтвердження
вимоги в межах заявленої моделі загроз і наявного пакета дослідницьких
матеріалів, а не абсолютне або повне доведення в усіх постановках.
Позначення «Підтримано інтерпретаційно» для R3 вказує на комбінатор,
сумісний з AND-інтерпретацією та підкріплений символьною моделлю, але
не на завершене обчислювальне доведення. Формулювання для R5 означає
практичну реалізованість і відтворюваність дослідницького матеріалу, а
не заявку на завершене промислове супроводження в усіх середовищах.

\subsection{Значення формальної верифікації}

Ключовим результатом є не «повне формальне доведення» всієї системи, а
узгодження кількох незалежних джерел доказовості. ProVerif краще
масштабується для аналізу секретності й кореспондентних властивостей у
необмеженій кількості сесій, тоді як Tamarin придатніший для
станових тверджень на кшталт прямої секретності та AND-моделі. Саме
поєднання цих інструментів із Rust-тестами дозволяє робити виваженіші
висновки, ніж опора на будь-яке одне джерело доказовості окремо.

\textbf{Лема P6 як центральний результат.} Вона показує, що у
сурогатній Tamarin-моделі компрометації довготривалих ключів
недостатньо для відновлення сесійного ключа без розкриття ефемерних
секретів. Це не замінює повного обчислювального доведення реалізації,
але становить сильну символьну аргументацію на користь того, що гібридний
комбінатор не зводиться до «класичного або постквантового шару окремо».

\subsection{Продуктивність та її інтерпретація}

Домінування Argon2id (порядку 510~мс в ізольованому вимірюванні) є
\textbf{свідомим архітектурним рішенням}: саме висока обчислювальна
вартість різко підвищує ціну офлайн-перебирання у сценаріях, де
противник не має окремого доступу до \texttt{oprf\_seed}. Постквантові операції
(122~мкс для ML-KEM-768, $\approx 164$~мкс для 4DH) на вимірювальній
платформі є малим доповненням порівняно з Argon2id.

Серверна пропускна здатність потребує акуратної інтерпретації. Показник
1,31~млн/с стосується лише фінальної функції \texttt{responder\_finish},
тоді як повний серверний цикл автентифікації обмежується фазою
\texttt{generate\_ke2} і становить приблизно 3,9~тис. автентифікацій на
секунду на ядро. Для багатьох інтерактивних сценаріїв цей рівень
пропускної здатності виглядає практично придатним, хоча остаточний
висновок залежить від профілю навантаження, доступу до сховища,
мережевої латентності та політик обмеження частоти
запитів.

З погляду сприйманої користувачем затримки це означає, що перехід до
Hybrid PQ-OPAQUE не змінює суттєво часовий профіль автентифікації.
Клієнтська сторона, як і раніше, здебільшого очікує завершення
Argon2id, тоді як сервер виконує криптографічне ядро менш ніж за
мілісекунду. Отже, питання продуктивності в практичному впровадженні
доцільніше формулювати не як ``чи є ML-KEM-768 надто дорогим?'', а як
``який профіль Argon2id є прийнятним для конкретного класу користувачів
і пристроїв?''.

\subsection{Стратегія міграції та зворотна сумісність}

Практичне впровадження Hybrid PQ-OPAQUE в інфраструктуру, де вже
використовується класичний OPAQUE або OPAQUE-подібний робочий процес,
потребує керованої процедури переходу. Прикладна перевага поточної
конструкції полягає в тому, що реєстраційний запис не збільшується:
сервер і надалі зберігає лише конверт і відкритий ключ клієнта, тоді як
постквантові елементи є повністю ефемерними і генеруються на кожну
сесію окремо. Це означає, що міграція майже не зачіпає схему бази
даних і здебільшого концентрується на прикладній логіці автентифікації.

Раціональна процедура переходу складається з трьох фаз.
\begin{enumerate}
  \item \textbf{Оновлення серверної логіки.} Сервер навчається приймати
        нові повідомлення протоколу, перевіряти 1-байтовий префікс версії та
        коректно обробляти \texttt{account\_id} як криптографічно
        зв'язаний контекст.
  \item \textbf{Поступова повторна реєстрація клієнтів.} Після
        оновлення клієнтських застосунків користувачі проходять
        повторну реєстрацію або перевипуск конверта в межах штатного
        життєвого циклу облікового запису, без масового перетворення
        вже збережених рядків БД.
  \item \textbf{Виведення класичного режиму з експлуатації.} Після
        завершення перехідного вікна сервер відхиляє старий формат
        повідомлень. Така навмисна несумісність потрібна для захисту від
        downgrade-атак і є частиною безпекового дизайну, а не
        експлуатаційним недоліком.
\end{enumerate}

Узгодження версії відбувається не лише на рівні парсингу, а й усередині
транскрипту через доменно розділені мітки та контекст
\texttt{account\_id}. Тому противник не може безкарно ``переключити''
сторону на інший режим роботи, залишивши решту криптографічного
контексту незмінною: будь-яка невідповідність змінює транскриптний геш
і призводить до розходження похідних MAC.

\subsection{Параметри Argon2id як головний важіль безпеки}

Поточна реалізація використовує конкретний профіль Argon2id:
$t = 3$ проходи, $m = 256$~MiB пам'яті, $p = 1$ смуга паралелізму та
16-байтова сіль. Саме цей профіль, а не гібридний KEM-компонент,
визначає основну частину наслідків для латентності й вартості
офлайн-перебирання~\cite{biryukov2016}. На вимірювальній платформі
роботи одна спроба обчислення похідного парольного значення коштує
\[
  C_{\mathrm{guess}}^{\mathrm{M1Pro}} \approx 0{,}51086~\text{с},
  \qquad
  R_{\mathrm{guess}}^{(1~core)} \approx 1{,}96~\text{спроб/с}.
\]
Ці величини не є універсальними константами противника; вони слугують
відтворюваним орієнтиром для порівняння внутрішніх накладних витрат
самого протоколу.

Важливішою за абсолютний час є властивість пам'яттєвої жорсткості цієї функції. При
$m = 256$~MiB масове перебиральне середовище швидко впирається у
дефіцит пам'яті, а не лише в арифметичну продуктивність. Саме тому
зменшення параметра пам'яті на практиці значно небезпечніше, ніж
оптимізація десятків мікросекунд у класичному чи постквантовому шарі
AKE. Отже, гібридизація підсилює захист від квантового противника, але
\emph{не знімає} відповідальності за правильне налаштування KSF.

\begin{table}[htbp]
\caption{Класи налаштування Argon2id для різних сценаріїв}
\label{tab:argon2-profiles}
\centering\small
\begin{tabularx}{\linewidth}{L{3.0cm}L{2.7cm}L{2.4cm}L{2.8cm}X}
\toprule
Клас & Приклад параметрів & Латентність & Сценарій & Інтерпретація \\
\midrule
Поточна реалізація & $t=3$, $m=256$~MiB, $p=1$ & 510{,}86~мс (виміряно) & Серверна автентифікація загального призначення & Базовий профіль роботи; максимальна частка словникової стійкості за поточного дизайну \\
Чутливий до латентності & $t=1$, $m=64$~MiB, $p=1$ & Менше поточного профілю & Мобільні клієнти, повільна пам'ять, сценарії з критичним користувацьким досвідом & Допустимий лише разом із жорстким обмеженням частоти запитів, моніторингом та додатковими антиавтоматизаційними засобами \\
Підвищена стійкість & $t\ge 3$, $m\ge 256$~MiB & Більше поточного профілю & Адміністративні, B2B або високоцінні акаунти & Підсилює словникову стійкість ціною більшої латентності; потребує окремого тестування доступності \\
\bottomrule
\end{tabularx}
\end{table}

Другий і третій рядки таблиці~\ref{tab:argon2-profiles} є не окремо
виміряними конфігураціями цієї роботи, а проектними орієнтирами.
Їх наведено навмисно, оскільки практична придатність парольного
протоколу визначається не лише криптографічною схемою, а й тим, як
оператор системи обирає положення на компромісній залежності
``словникова стійкість~--- латентність''.

Характерно, що постквантові накладні витрати $\Delta_{PQ} = 159{,}3$~мкс
залишаються на кілька порядків меншими за будь-який реалістичний
профіль Argon2id. Тому з практичного погляду відповідь на питання
``чи можна дозволити собі ML-KEM-768 у PAKE?'' є позитивною; справжній
операційний вибір стосується параметрів KSF, а не самого факту
постквантової гібридизації.

\subsection{Операційні припущення та межа довіри}

Сильні безпекові твердження для асиметричного PAKE коректні лише тоді, коли
чітко окреслено, який саме серверний секрет вважається компрометованим.
У цій реалізації база даних із реєстраційними записами та глобальне
\texttt{oprf\_seed} утворюють різні рівні довіри. Компрометація лише
пасивного сховища не надає противнику готового парольного еквівалента;
натомість компрометація \texttt{oprf\_seed} дозволяє детерміновано
відтворювати OPRF-ключі для окремих облікових записів та перетворює викрадені записи на
матеріал для офлайн-перевірки словникових кандидатів. Саме тому
твердження про словникову стійкість у статті всюди прив'язуються до
сценарію компрометації бази даних \emph{без} компрометації
\texttt{oprf\_seed}.

З експлуатаційного погляду це означає, що \texttt{oprf\_seed} слід
захищати не як ``допоміжну конфігурацію'', а як кореневий
криптографічний секрет сервісу. Практичними наслідками є використання
окремого сховища секретів, рознесення доступів між БД і KMS/HSM,
мінімізація часу присутності насіння в пам'яті процесу та ретельна
політика ротації. Контекстне зв'язування \texttt{account\_id} також
набуває операційного значення: помилка в нормалізації ідентифікаторів
не обов'язково зламає примітиви безпосередньо, але здатна створити
небезпечну семантичну неоднозначність між обліковими записами.

Отже, наукову новизну гібридної схеми потрібно читати разом із
реалістичною моделлю розгортання. Символьна лема про словникову
стійкість, тести на граничні умови та інженерна політика поводження з
\texttt{oprf\_seed} слід розглядати як взаємодоповнювальні елементи
єдиної аргументації безпеки.

\subsection{Транспортні та мережеві наслідки формату передавання повідомлень}

Запропонована конструкція зберігає кількість раундів обміну, однак
збільшує обсяг повідомлень. У цьому підрозділі ці наслідки
інтерпретуються саме в термінах мережевого формату повідомлень,
визначеного в Табл.~\ref{tab:terminology}. Для транспорту поверх TCP
або QUIC це здебільшого означає лише більший обсяг переданих байтів і вищі вимоги
до буферів. Для датаграм-орієнтованих середовищ наслідки помітніші:
повідомлення KE1 (1273~байти у мережевому форматі) і KE2 (1377~байтів у мережевому форматі) можуть
перетинати практичну межу безпечного обсягу корисних даних для окремої UDP-датаграми,
особливо в мережах з мінімальним MTU IPv6.

З цієї причини формат передавання повідомлень не може розглядатися як другорядна деталь.
У системах із власним кадруванням поверх UDP фрагментація збільшує
ймовірність втрати пакета, латентність повторної передачі та складність
моніторингу. Натомість у TCP-/QUIC-сценаріях збільшення обсягу не змінює
кількості логічних раундів і тому не погіршує асимптотики протоколу.
За таких умов Hybrid PQ-OPAQUE залишається придатним для
інтерактивної автентифікації, але середовище транспорту починає
відігравати істотно більшу роль, ніж у класичному OPAQUE з коротшими
повідомленнями.

Позитивним наслідком точного опису мережевого формату є те, що оператор системи може
планувати мережевий бюджет до впровадження протоколу: оцінювати MTU,
кадрування, ліміти API-шлюзів і поведінку WAF/IDS щодо великих бінарних
полів. Саме така можливість раннього планування і є одним із
методологічних аргументів на користь детального реалізаційного опису в
науковій статті, а не лише наведення логічної схеми обміну.

\subsection{Сценарії розгортання та масштабування}

Криптографічна придатність протоколу ще не гарантує його однакової
експлуатаційної придатності в різних класах систем. Для
Hybrid PQ-OPAQUE визначальним є нерівномірний розподіл джерел затримки:
на клієнті домінує Argon2id, на сервері
--- генерація KE2, а на мережі --- транспортні властивості великих
повідомлень у мережевому форматі.
Тому одна й та сама конструкція може бути оптимальною для
централізованого веб-сервісу, прийнятною для мобільного застосунку та
надмірно вимогливою для крайового вузла з датаграмним транспортом.

\begin{table}[htbp]
\caption{Типові сценарії розгортання протоколу Hybrid PQ-OPAQUE}
\label{tab:deployment-scenarios}
\centering\small
\begin{tabularx}{\linewidth}{L{3.0cm}L{3.2cm}L{2.6cm}X}
\toprule
Сценарій & Домінуюче вузьке місце & Рекомендований клас KSF & Практичний висновок \\
\midrule
Веб/SaaS автентифікація & Argon2id на клієнті; введення-виведення на сервері & Поточна реалізація & Найсприятливіший сценарій: постквантовий шар не змінює істотно користувацький досвід, а серверний обчислювальний резерв залишається значним \\
Мобільний клієнт & Пропускна здатність пам'яті пристрою & Чутливий до латентності & Доцільний окремий профіль KSF і контроль зловживань; при цьому KEM-накладні витрати залишаються другорядними \\
Адміністративний/привілейований доступ & Стійкість до офлайн-перебирання & Підвищена стійкість & Доцільним є прийняття більшої латентності заради посилення KSF, оскільки частота входів низька, а цінність облікових записів висока \\
Датаграмне/крайове середовище & MTU, фрагментація, втрата пакетів & Залежить від платформи & Потрібно окремо верифікувати транспортний бюджет; розмір повідомлень у мережевому форматі стає не менш важливим, ніж криптографічний час \\
\bottomrule
\end{tabularx}
\end{table}

Особливо важливою є теоретична межа серверної пропускної здатності. Для фази респондера
емпірично маємо
\[
  Q_{\mathrm{core}} \approx
  \frac{1}{254{,}26~\mu s + 0{,}76082~\mu s}
  \approx 3{,}92 \times 10^3
  \text{ автентифікацій/с на ядро}.
\]
Це значення слід читати як \emph{оптимістичну верхню межу} для
криптографічно обмеженого сервера без урахування доступу до БД,
мережевої черги, телеметрії, бізнес-логіки та політик обмеження
частоти запитів. Проте навіть така стримана оцінка означає, що
8-ядерний вузол має криптографічний резерв порядку десятків тисяч
автентифікацій за секунду, тобто для більшості веб- і B2B-сценаріїв
вузьким місцем буде не KE2, а периферійні компоненти системи.

З наукового погляду цей висновок важливий тому, що він відокремлює
``криптографічну вартість'' від ``системної вартості''. Стаття показує,
що постквантовий захист на рівні AKE не руйнує масштабованість сам по
собі; складність міграції та експлуатації переноситься в домен
налаштування KSF, транспортного кадрування та захисту кореневих
секретів. Саме така декомпозиція і робить запропонований протокол
реалістичним кандидатом для впровадження, а не лише лабораторною
конструкцією.

\subsection{Відтворюваність і валідність результатів}

Наукова цінність системної роботи визначається не лише теоретичною
новизною, а й тим, наскільки результати може незалежно перевірити інший
дослідник. Істотною перевагою поточної статті є наявність
відкритого пакета дослідницьких матеріалів: вихідних кодів, формальних моделей,
архівованих логів Tamarin/ProVerif, тестового набору й модулів
Criterion. Це дає змогу відокремити власне наукове твердження від
одноразового «авторського прогону».

Водночас відтворюваність не є повною в строгому сенсі. Часи
еталонних вимірювань отримано на одній апаратній платформі, а сторінка
результатів містить числа, чутливі до версій компілятора, бібліотек і
системного планувальника. Тому в роботі доцільно інтерпретувати
абсолютні мікросекундні значення як \emph{репрезентативні} для
порівняння відносних порядків величин, а не як універсальні константи
для всіх середовищ розгортання.

\subsection{Поточні обмеження}

\begin{enumerate}
  \item \textbf{Відповідність IETF.} Реалізація є «OPAQUE-подібною»
        і не є повністю конформною до
        RFC~9807 та пов'язаних стандартних інтерфейсів OPAQUE~\cite{ietf-opaque}.
  \item \textbf{Межа довіри до серверних секретів.} Твердження про
        словникову стійкість коректні лише для сценарію компрометації
        бази даних без компрометації \texttt{oprf\_seed}; це
        фундаментальна межа моделі, а не другорядна примітка.
  \item \textbf{Рівень зрілості постквантових стандартів.} ML-KEM-768
        (FIPS~203)~--- відносно новий стандарт (серпень~2024).
        Гібридизація зменшує залежність від стійкості лише одного
        криптографічного припущення, але не усуває потреби в подальшому
        аналізі стандарту.
  \item \textbf{Символьна та обчислювальна моделі.} Tamarin/ProVerif
        верифікують символьну модель~\cite{cremers2012}. Розрив між
        рівнями є стандартним обмеженням формальної
        верифікації~\cite{basin2018}.
  \item \textbf{Збільшений обсяг повідомлень у мережевому форматі.} Повідомлення
        KE1/KE2 істотно більші за класичний OPAQUE, що може створювати
        додаткові мережеві накладні витрати у датаграм-орієнтованих
        середовищах.
  \item \textbf{Одна апаратна платформа вимірювання.} Абсолютні
        значення латентності отримано на Apple~M1~Pro; переносити їх на
        інші архітектури слід обережно.
\end{enumerate}

\subsection{Напрями подальших досліджень}

\begin{enumerate}
  \item \textbf{Конформність IETF.} Адаптація до
        RFC~9807 та пов'язаної екосистеми OPAQUE/OPRF.
  \item \textbf{Постквантові цифрові підписи.} Дослідження інтеграції
        ML-DSA (FIPS~204) для підписування зовнішніх протокольних
        повідомлень або метаданих у сценаріях, де потрібна
        невідмовність; це не замінює HMAC у базовому рукостисканні.
  \item \textbf{ML-KEM-512 для пристроїв IoT.} Оцінка компромісу
        безпека/ефективність.
  \item \textbf{UC-модель.} Формальне доведення в рамках Universal
        Composability~\cite{canetti2001} для постквантових розширень.
  \item \textbf{Стандартизація IETF CFRG.} Розробка пропозиції
        щодо стандартизованого постквантового розширення OPAQUE.
\end{enumerate}

% ============================================================
\section{Висновки}
% ============================================================

Отримані результати дають підстави розглядати гібридизацію OPAQUE з
ML-KEM-768 як цілісну науково-інженерну конструкцію з формально
описаним форматом передавання повідомлень, відтворюваними моделями,
повноцінною Rust-реалізацією та кількісною оцінкою накладних витрат.

\textbf{Теоретичний внесок.} Запропоновано конструкцію Hybrid
PQ-OPAQUE, у якій ML-KEM-768 інтегровано як невід'ємну складову
протоколу. Ключовий матеріал формується комбінатором
\begin{equation}
  \mathrm{PRK} = \mathrm{HKDF\text{-}Extract}\!\bigl(
    \mathrm{salt}_{\tau},\;
    dh_1 \| dh_2 \| dh_3 \| dh_4 \| ss_{\mathrm{kem}}\bigr)
\end{equation}
з розширеним транскриптом; у межах прийнятої моделі він є сумісним з
AND-інтерпретацією безпеки на рівні ключового матеріалу
(див. інтерпретаційне твердження в розд.~4).
Четвертий добуток Діффі--Геллмана ($dh_4 = e_S \cdot E_C$)
посилює зв'язування сторін; його можливий внесок у стійкість до
UKS-подібних підстановок у цій роботі розглядається як дизайн-мотивація.
Успадковані від OPAQUE цілі aPAKE у цій роботі розглядаються крізь
сукупність тверджень про межі офлайн-перебирання, тверджень про пряму
секретність і взаємну автентифікацію, а також символьних моделей і
тестових матеріалів у межах явно зафіксованої моделі загроз.

\textbf{Верифікаційний внесок.}
\begin{itemize}
  \item \textbf{ProVerif~2.05}: 5/5 запитів підтверджено у двох
        окремих моделях (секретність та автентифікація);
  \item \textbf{Tamarin~Prover~1.10.0}: 8/8 лем підтверджено в
        архівованій сурогатній 4DH-моделі за 22,83~с; зокрема
        пряма секретність (119~кроків) та AND-модель (29~кроків);
  \item формальні моделі та журнали верифікації інтерпретуються як
        символьна доказова база,
        а не як машинно перевірене доведення коду реалізації.
\end{itemize}

\textbf{Реалізаційний та вимірювальний внесок.}
\begin{itemize}
  \item \textbf{155 тестів пройдено, 1 пропущено} у поточному просторі проєктів;
  \item $\texttt{authentication\_e2e}$ = \textbf{481,43~мс}
        (95\%~ДІ [471,73; 495,25]); ізольований профіль Argon2id =
        510,86~мс; ML-KEM-768 = 122,01~мкс. Ці числа походять з окремих сценаріїв
        Criterion і не є адитивними;
  \item ізольована серверна фаза \texttt{ke2\_only} = \textbf{254,26~мкс},
        тоді як \texttt{responder\_finish} = 760,82~нс;
  \item постквантові накладні витрати: \textbf{$+159{,}3$~мкс} ($+30{,}1\%$)
        без KSF; \textbf{близько 0,03\%} від наскрізного часу;
        $+2272$~байти у мережевому форматі ($+516\%$) при незмінних 1,5~обмінах.
\end{itemize}

У сукупності ці результати показують, що Hybrid PQ-OPAQUE слід
розглядати не як абстрактну схему, а як відтворюваний крок до
постквантового підсилення парольної автентифікації. Принципово
важливо, що додавання ML-KEM-768 не змінює кількості раундів і не
збільшує суттєво обчислювальну вартість відносно KSF; основний
компроміс зміщується в мережевий обсяг рукостискання. У межах
прийнятої моделі гібридний комбінатор зберігає змістовну перехідну
цінність доти, доки принаймні одне з двох джерел ключового матеріалу
лишається обчислювально стійким. Водночас практичне читання цих
результатів потребує академічної стриманості: формальні висновки слід
інтерпретувати разом із межами моделі, а твердження про словникову
стійкість~--- лише для сценарію компрометації бази даних без
компрометації \texttt{oprf\_seed}. Головна цінність роботи полягає
саме в поєднанні формальної моделі, відкритої реалізації та
відтворюваного кількісного результату.

% ============================================================
% Канонічну бібліографію для конкурсного PDF підтримуємо безпосередньо в
% main.tex, щоб документ залишався самодостатнім без зовнішнього BibTeX-кроку.
\begin{thebibliography}{34}
% ============================================================

\bibitem{bonneau2012}
Bonneau~J., Herley~C., van~Oorschot~P.\,C., Stajano~F.
The quest to replace passwords: A framework for comparative evaluation of web authentication schemes~//
Proc. IEEE S\&P~2012. --- P.~553--567.

\bibitem{florencio2007}
Florencio~D., Herley~C.
A large-scale study of web password habits~//
Proc. WWW~2007. --- P.~657--666.

\bibitem{bellare2000}
Bellare~M., Pointcheval~D., Rogaway~P.
Authenticated key exchange secure against dictionary attacks~//
EUROCRYPT~2000, LNCS~1807. --- Springer, 2000. --- P.~139--155.

\bibitem{jarecki2018}
Jarecki~S., Krawczyk~H., Xu~J.
OPAQUE: An asymmetric PAKE protocol secure against pre-computation attacks~//
EUROCRYPT~2018, LNCS~10822. --- Springer, 2018. --- P.~456--486.

\bibitem{ietf-opaque}
Bourdrez~D., Krawczyk~H., Lewi~K., Wood~C.\,A.
The OPAQUE Augmented Password-Authenticated Key Exchange (aPAKE) Protocol~//
IRTF RFC~9807. --- July~2025. ---
URL:~\url{https://www.rfc-editor.org/info/rfc9807}.

\bibitem{shor1994}
Shor~P.\,W.
Algorithms for quantum computation: Discrete logarithms and factoring~//
FOCS~1994. --- P.~124--134.

\bibitem{bernstein2017}
Bernstein~D.\,J., Lange~T.
Post-quantum cryptography~//
Nature. --- 2017. --- Vol.~549, No.~7671. --- P.~188--194.

\bibitem{roetteler2017}
Roetteler~M., Naehrig~M., Svore~K.\,M., Lauter~K.
Quantum resource estimates for computing elliptic curve discrete logarithms~//
ASIACRYPT~2017, LNCS~10625. --- Springer, 2017. --- P.~241--270.

\bibitem{babbush2026}
Babbush~R., Zalcman~A., Gidney~C., Broughton~M., Khattar~T.,
Neven~H., Bergamaschi~T., Drake~J., Boneh~D.
Securing Elliptic Curve Cryptocurrencies against Quantum Vulnerabilities:
Resource Estimates and Mitigations~//
Google Quantum AI Whitepaper. --- March~30,~2026. ---
URL:~\url{https://quantumai.google/static/site-assets/downloads/cryptocurrency-whitepaper.pdf}.

\bibitem{mosca2018}
Mosca~M.
Cybersecurity in an era with quantum computers: Will we be ready?~//
IEEE Security \& Privacy. --- 2018. --- Vol.~16, No.~5. --- P.~38--41.

\bibitem{fips203}
National Institute of Standards and Technology.
Module-Lattice-Based Key-Encapsulation Mechanism Standard: FIPS~203. --- 2024. ---
URL:~\url{https://csrc.nist.gov/pubs/fips/203/final}.

\bibitem{fips204}
National Institute of Standards and Technology.
Module-Lattice-Based Digital Signature Standard: FIPS~204. --- 2024. ---
URL:~\url{https://csrc.nist.gov/pubs/fips/204/final}.

\bibitem{stebila2024}
Stebila~D., Fluhrer~S., Gueron~S.
Hybrid key exchange in TLS~1.3~//
IETF Internet-Draft draft-ietf-tls-hybrid-design-16 (work in progress). ---
September~2025. ---
URL:~\url{https://datatracker.ietf.org/doc/draft-ietf-tls-hybrid-design/}.

\bibitem{brendel2020}
Brendel~J., Fiedler~R., G\"{u}nther~F., Jacobsen~C., Poettering~B.
Post-quantum security of the Signal protocol~//
SAC~2020, LNCS~12804. --- Springer, 2020. --- P.~567--591.

\bibitem{hulsing2021}
H\"{u}lsing~A., Ning~K.-C., Schwabe~P., Weber~F., Zimmermann~P.\,R.
Post-quantum WireGuard~//
IEEE S\&P~2021. --- P.~304--321.

\bibitem{boyko2000}
Boyko~V., MacKenzie~P., Patel~S.
Provably secure password-authenticated key exchange using Diffie-Hellman~//
EUROCRYPT~2000, LNCS~1807. --- Springer, 2000. --- P.~156--171.

\bibitem{bellovin1992}
Bellovin~S.\,M., Merritt~M.
Encrypted key exchange: Password-based protocols secure against dictionary attacks~//
IEEE S\&P~1992. --- P.~72--84.

\bibitem{ladd2023}
Ladd~W., Kaduk~B., Harkins~D.
SPAKE2, a PAKE~// IETF RFC~9382. --- 2023.

\bibitem{haase2019}
Haase~B., Labrique~B.
AuCPace: Efficient verifier-based PAKE protocol tailored for the IIoT~//
IACR TCHES. --- 2019. --- Vol.~2019, No.~2. --- P.~1--48.

\bibitem{wu1998}
Wu~T.
The Secure Remote Password protocol~// NDSS~1998. --- P.~97--111.

\bibitem{devalence2024}
de~Valence~H., Grigg~J., Hamburg~M., Lovecruft~I., Tankersley~G.,
Valsorda~F.
The ristretto255 and decaf448 groups~//
IRTF RFC~9496. --- December~2023. ---
URL:~\url{https://www.rfc-editor.org/info/rfc9496}.

\bibitem{canetti2001}
Canetti~R.
Universally composable security: A new paradigm for cryptographic protocols~//
FOCS~2001. --- P.~136--145.

\bibitem{krawczyk2010}
Krawczyk~H., Eronen~P.
HMAC-based Extract-and-Expand Key Derivation Function (HKDF)~//
IETF RFC~5869. --- 2010.

\bibitem{cremers2012}
Cremers~C., Mauw~S.
Operational Semantics and Verification of Security Protocols. --- Springer, 2012.

\bibitem{basin2018}
Basin~D., Cremers~C., Meadows~C.
Model checking security protocols~//
Handbook of Model Checking. --- Springer, 2018. --- P.~727--762.

\bibitem{biryukov2016}
Biryukov~A., Dinu~D., Khovratovich~D.
Argon2: New generation of memory-hard functions for password hashing~//
IEEE Euro S\&P~2016. --- P.~292--302.

\bibitem{heyer2023}
Heyer~C. et al.
Criterion.rs: Statistics-driven micro-benchmarking in Rust. ---
URL:~\url{https://github.com/bheisler/criterion.rs}. --- 2023.

\bibitem{beguinet2023}
Beguinet~H. et al.
Get Me out of This Lattice: CRYSTALS-Kyber and Post-Quantum Signal~//
ACM CCS~2023.

\bibitem{arriaga2024}
Arriaga~A., Barbosa~M., Jarecki~S., Skrobot~M.
C'est tr\`es CHIC: A compact password-authenticated key exchange from lattice-based KEM~//
Cryptology ePrint Archive. --- Paper~2024/308. --- 2024. ---
URL:~\url{https://eprint.iacr.org/2024/308}.

\bibitem{hesse2024}
Hesse~J., Rosenberg~M.
PAKE combiners and efficient post-quantum instantiations~//
EUROCRYPT~2025, LNCS~15602. --- Springer, 2025. --- P.~395--420.

\bibitem{lyu2024uc}
Lyu~Y., Liu~S., Han~S.
Universal composable password authenticated key exchange for the post-quantum world~//
EUROCRYPT~2024, LNCS~14657. --- Springer, 2024. --- P.~120--150.

\bibitem{lyu2024hybrid}
Lyu~Y., Liu~S.
Hybrid password authentication key exchange in the UC framework~//
EUROCRYPT~2025, LNCS~15602. --- Springer, 2025. --- P.~421--450.

\bibitem{vos2025}
Vos~J., Jarecki~S., Wood~C.\,A.
Hybrid Post-Quantum Password Authenticated Key Exchange~//
Individual Internet-Draft draft-vos-cfrg-pqpake-00. --- April~2025. ---
URL:~\url{https://datatracker.ietf.org/doc/draft-vos-cfrg-pqpake/00/}.

\bibitem{castryck2023}
Castryck~W., Decru~T.
An efficient key recovery attack on SIDH~//
EUROCRYPT~2023, LNCS~14008. --- Springer, 2023. --- P.~423--447.

\end{thebibliography}

\end{document}
